%% Generated by tools/gen_error_codes.py from docs/errors/ of the compiler;
%% edit the compiler's pages or tools/error_themes.txt, then rerun it.

\clearpage
\section{Classes, records and traits}
\label{sec:codes:classes}

\begin{longtable}{@{}l p{0.78\linewidth} r@{}}
\toprule Code & Message & Page\\ \midrule \endhead
\hyperref[sec:codes:e3001]{E3001} & a class cannot be both abstract and final & \pageref{sec:codes:e3001}\\
\hyperref[sec:codes:e3006]{E3006} & the ctor built from the fields is already declared & \pageref{sec:codes:e3006}\\
\hyperref[sec:codes:e3007]{E3007} & the ctor built from the fields is implicit in a data record & \pageref{sec:codes:e3007}\\
\hyperref[sec:codes:e3010]{E3010} & a class may only define one destructor & \pageref{sec:codes:e3010}\\
\hyperref[sec:codes:e3015]{E3015} & a record cannot be both overlaid and packed & \pageref{sec:codes:e3015}\\
\hyperref[sec:codes:e3030]{E3030} & unexpected declaration inside an impl block & \pageref{sec:codes:e3030}\\
\hyperref[sec:codes:e3033]{E3033} & unexpected declaration inside a trait & \pageref{sec:codes:e3033}\\
\hyperref[sec:codes:e4011]{E4011} & cannot cast derived type \errhole{} to base class type \errhole{} in that context & \pageref{sec:codes:e4011}\\
\hyperref[sec:codes:e4017]{E4017} & field \errhole{} has not yet been validated & \pageref{sec:codes:e4017}\\
\hyperref[sec:codes:e4018]{E4018} & class type \errhole{} does not implement trait \errhole{} & \pageref{sec:codes:e4018}\\
\hyperref[sec:codes:e4019]{E4019} & class \errhole{} has no field named \errhole{} & \pageref{sec:codes:e4019}\\
\hyperref[sec:codes:e4020]{E4020} & class \errhole{} has no method named \errhole{} & \pageref{sec:codes:e4020}\\
\hyperref[sec:codes:e4024]{E4024} & colliding method definitions \errhole{} and \errhole{} & \pageref{sec:codes:e4024}\\
\hyperref[sec:codes:e4026]{E4026} & cannot construct an instance of class \errhole{} that is abstract & \pageref{sec:codes:e4026}\\
\hyperref[sec:codes:e4027]{E4027} & method \errhole{} is defined as constant & \pageref{sec:codes:e4027}\\
\hyperref[sec:codes:e4035]{E4035} & infinite constructor redirection calls when calling \errhole{} & \pageref{sec:codes:e4035}\\
\hyperref[sec:codes:e4044]{E4044} & field must have a referenceable name & \pageref{sec:codes:e4044}\\
\hyperref[sec:codes:e4051]{E4051} & cannot perform a direct call of the method \errhole{} that is empty & \pageref{sec:codes:e4051}\\
\hyperref[sec:codes:e4060]{E4060} & field access method must return a value & \pageref{sec:codes:e4060}\\
\hyperref[sec:codes:e4062]{E4062} & field assign method must be mutable & \pageref{sec:codes:e4062}\\
\hyperref[sec:codes:e4063]{E4063} & field assign method must not return a value & \pageref{sec:codes:e4063}\\
\hyperref[sec:codes:e4065]{E4065} & field method \errhole{} is mutable & \pageref{sec:codes:e4065}\\
\hyperref[sec:codes:e4066]{E4066} & field method cannot have more than one parameter & \pageref{sec:codes:e4066}\\
\hyperref[sec:codes:e4079]{E4079} & the field \errhole{} is not public, and has no initial value to be constructed from & \pageref{sec:codes:e4079}\\
\hyperref[sec:codes:e4090]{E4090} & method \errhole{} implicitly overrides method \errhole{} & \pageref{sec:codes:e4090}\\
\hyperref[sec:codes:e4091]{E4091} & implicit override of method \errhole{} by trait impl & \pageref{sec:codes:e4091}\\
\hyperref[sec:codes:e4093]{E4093} & trait \errhole{} is implemented multiple times by \errhole{} & \pageref{sec:codes:e4093}\\
\hyperref[sec:codes:e4094]{E4094} & impl statement must be followed by a trait, not \errhole{} & \pageref{sec:codes:e4094}\\
\hyperref[sec:codes:e4101]{E4101} & the base class \errhole{} is marked as final & \pageref{sec:codes:e4101}\\
\hyperref[sec:codes:e4102]{E4102} & the base of a class must be a class, not a \errhole{} & \pageref{sec:codes:e4102}\\
\hyperref[sec:codes:e4103]{E4103} & method \errhole{} cannot be inline as it is virtual(overridable or inherited) & \pageref{sec:codes:e4103}\\
\hyperref[sec:codes:e4133]{E4133} & field \errhole{} is initialized multiple times & \pageref{sec:codes:e4133}\\
\hyperref[sec:codes:e4141]{E4141} & method \errhole{} is defined as mutable & \pageref{sec:codes:e4141}\\
\hyperref[sec:codes:e4144]{E4144} & class \errhole{} is not abstract but has empty methods & \pageref{sec:codes:e4144}\\
\hyperref[sec:codes:e4147]{E4147} & \errhole{} is not an ancestor type of \errhole{} & \pageref{sec:codes:e4147}\\
\hyperref[sec:codes:e4150]{E4150} & \errhole{} is not a class type & \pageref{sec:codes:e4150}\\
\hyperref[sec:codes:e4162]{E4162} & \errhole{} is not a trait & \pageref{sec:codes:e4162}\\
\hyperref[sec:codes:e4168]{E4168} & no constructor found for class \errhole{} & \pageref{sec:codes:e4168}\\
\hyperref[sec:codes:e4169]{E4169} & no constructor named \errhole{} found for class \errhole{} & \pageref{sec:codes:e4169}\\
\hyperref[sec:codes:e4172]{E4172} & record or entity type \errhole{} has no size due to forward references & \pageref{sec:codes:e4172}\\
\hyperref[sec:codes:e4173]{E4173} & class \errhole{} has no ancestor & \pageref{sec:codes:e4173}\\
\hyperref[sec:codes:e4181]{E4181} & method \errhole{} must have a body to override \errhole{} & \pageref{sec:codes:e4181}\\
\hyperref[sec:codes:e4182]{E4182} & method \errhole{} must be a field method to override \errhole{} & \pageref{sec:codes:e4182}\\
\hyperref[sec:codes:e4183]{E4183} & cannot override final method \errhole{} & \pageref{sec:codes:e4183}\\
\hyperref[sec:codes:e4184]{E4184} & the return type of the overriding method \errhole{} is not compatible with the return type of the ancestor method \errhole{} & \pageref{sec:codes:e4184}\\
\hyperref[sec:codes:e4185]{E4185} & the protection \errhole{} of the overriding method \errhole{} does not match the definition in the ancestor class \errhole{} & \pageref{sec:codes:e4185}\\
\hyperref[sec:codes:e4186]{E4186} & the throwers of the overriding method \errhole{} are not compatible with the throwers of the ancestor method \errhole{} & \pageref{sec:codes:e4186}\\
\hyperref[sec:codes:e4187]{E4187} & trait method \errhole{} was already overridden & \pageref{sec:codes:e4187}\\
\hyperref[sec:codes:e4188]{E4188} & cannot override a non trait method \errhole{} with \errhole{} inside impl & \pageref{sec:codes:e4188}\\
\hyperref[sec:codes:e4189]{E4189} & method \errhole{} overrides nothing & \pageref{sec:codes:e4189}\\
\hyperref[sec:codes:e4190]{E4190} & field method \errhole{} cannot override the non field method \errhole{} & \pageref{sec:codes:e4190}\\
\hyperref[sec:codes:e4192]{E4192} & cannot override trait method \errhole{} with \errhole{} outside impl & \pageref{sec:codes:e4192}\\
\hyperref[sec:codes:e4196]{E4196} & ancestor cycle found when validating \errhole{} & \pageref{sec:codes:e4196}\\
\hyperref[sec:codes:e4233]{E4233} & no constructor is callable with the parameters \errhole{} & \pageref{sec:codes:e4233}\\
\hyperref[sec:codes:e4253]{E4253} & the field \errhole{} has no initial value & \pageref{sec:codes:e4253}\\
\hyperref[sec:codes:e4319]{E4319} & method \errhole{} must be a generator to override the generator \errhole{} & \pageref{sec:codes:e4319}\\
\hyperref[sec:codes:e4320]{E4320} & generator \errhole{} cannot override the non generator method \errhole{} & \pageref{sec:codes:e4320}\\
\bottomrule
\end{longtable}

\errorcode{E3001}{a class cannot be both abstract and final}
\label{sec:codes:e3001}

Raised during \textbf{declaration}, from \texttt{DeclareErrorMessage::\allowbreak{}ABSTRACT\_\allowbreak{}AND\_\allowbreak{}FINAL} (\textit{src/\allowbreak{}ymirc/\allowbreak{}semantic/\allowbreak{}declarator/\allowbreak{}errors.yr}).

\subsubsection*{What it means}

A class is declared both \tokennolst{abstract} and \tokennolst{final}. \tokennolst{abstract} means the class cannot be instantiated and must be derived from; \tokennolst{final} means it cannot be derived from. Together they describe a class nothing can ever do anything with.

\subsubsection*{Example}

\textit{test\_\allowbreak{}resources/\allowbreak{}diagnostics/\allowbreak{}test4.yr}:

%% check: skip
\begin{lstlisting}[style=coloredverbatimError]
@abstract @final
class Foo {
    pub self () {}
}

fn main () {}
\end{lstlisting}

\begin{lstlisting}[style=bashVerb, breaklines=true, escapechar=~]
Error[E3001] : a class cannot be both abstract and final
 --> test_resources/diagnostics/test4.yr:(1,12)
 1  ~\lstbox{\textSFxi{}}~ @abstract @final
    ~\lstbox{\textSFv{}}~            ^^^^^
    ~\lstbox{\textSFxi{}}~ Note :
    ~\lstbox{\textSFxi{}}~  --> test_resources/diagnostics/test4.yr:(1,2)
    ~\lstbox{\textSFxi{}}~  1  ~\lstbox{\textSFxi{}}~ @abstract @final
    ~\lstbox{\textSFxi{}}~     ~\lstbox{\textSFv{}}~  ^^^^^^^^
    ~\lstbox{\textSFxi{}}~     ~\lstbox{\textSFii{}}~~\lstbox{\textSFx{}}~~\lstbox{\textSFx{}}~~\lstbox{\textSFx{}}~~\lstbox{\textSFx{}}~~\lstbox{\textSFx{}}~~\lstbox{\textSFx{}}~~\lstbox{\textSFx{}}~
    ~\lstbox{\textSFxi{}}~ Note :
    ~\lstbox{\textSFxi{}}~  --> test_resources/diagnostics/test4.yr:(2,1)
    ~\lstbox{\textSFxi{}}~  2  ~\lstbox{\textSFxi{}}~ class Foo {
    ~\lstbox{\textSFxi{}}~     ~\lstbox{\textSFv{}}~ ^^^^^
    ~\lstbox{\textSFxi{}}~     ~\lstbox{\textSFii{}}~~\lstbox{\textSFx{}}~~\lstbox{\textSFx{}}~~\lstbox{\textSFx{}}~~\lstbox{\textSFx{}}~~\lstbox{\textSFx{}}~~\lstbox{\textSFx{}}~~\lstbox{\textSFx{}}~
    ~\lstbox{\textSFii{}}~~\lstbox{\textSFx{}}~~\lstbox{\textSFx{}}~~\lstbox{\textSFx{}}~~\lstbox{\textSFx{}}~~\lstbox{\textSFx{}}~~\lstbox{\textSFvii{}}~~\lstbox{\textSFx{}}~
\end{lstlisting}

\subsubsection*{How to fix it}

Drop whichever of the two does not apply.

\errorcode{E3006}{the ctor built from the fields is already declared}
\label{sec:codes:e3006}

Raised during \textbf{declaration}, from \texttt{DeclareErrorMessage::\allowbreak{}DEFAULT\_\allowbreak{}CTOR\_\allowbreak{}ALREADY\_\allowbreak{}DECLARED} (\textit{src/\allowbreak{}ymirc/\allowbreak{}semantic/\allowbreak{}declarator/\allowbreak{}errors.yr}).

\subsubsection*{What it means}

The constructor built from the fields of the type is declared explicitly while the compiler also generates it. Both would have the same signature.

\subsubsection*{Example}

\textit{test\_\allowbreak{}resources/\allowbreak{}default\_\allowbreak{}ctor/\allowbreak{}test11.yr}:

%% check: skip
\begin{lstlisting}[style=coloredverbatimError]
// The request is a declaration of its own, so a class makes it once.
class Foo {
   pub let x: i32;

   pub self;
   prv self;
}

fn main () {
    let a = copy Foo (1);
    a.x;
}
\end{lstlisting}

\begin{lstlisting}[style=bashVerb, breaklines=true, escapechar=~]
Error[E3006] : the ctor built from the fields is already declared
 --> test_resources/default_ctor/test11.yr:(6,8)
 6  ~\lstbox{\textSFxi{}}~    prv self;
    ~\lstbox{\textSFv{}}~        ^^^^
    ~\lstbox{\textSFxi{}}~ Note : declared here
    ~\lstbox{\textSFxi{}}~  --> test_resources/default_ctor/test11.yr:(5,8)
    ~\lstbox{\textSFxi{}}~  5  ~\lstbox{\textSFxi{}}~    pub self;
    ~\lstbox{\textSFxi{}}~     ~\lstbox{\textSFv{}}~        ^^^^
    ~\lstbox{\textSFxi{}}~     ~\lstbox{\textSFii{}}~~\lstbox{\textSFx{}}~~\lstbox{\textSFx{}}~~\lstbox{\textSFx{}}~~\lstbox{\textSFx{}}~~\lstbox{\textSFx{}}~~\lstbox{\textSFx{}}~~\lstbox{\textSFx{}}~
    ~\lstbox{\textSFii{}}~~\lstbox{\textSFx{}}~~\lstbox{\textSFx{}}~~\lstbox{\textSFx{}}~~\lstbox{\textSFx{}}~~\lstbox{\textSFx{}}~~\lstbox{\textSFvii{}}~~\lstbox{\textSFx{}}~
\end{lstlisting}

\subsubsection*{How to fix it}

Remove the explicit constructor and rely on the generated one, or change its parameters so the two no longer collide.

\errorcode{E3007}{the ctor built from the fields is implicit in a data record}
\label{sec:codes:e3007}

Raised during \textbf{declaration}, from \texttt{DeclareErrorMessage::\allowbreak{}DEFAULT\_\allowbreak{}CTOR\_\allowbreak{}IN\_\allowbreak{}DATA} (\textit{src/\allowbreak{}ymirc/\allowbreak{}semantic/\allowbreak{}declarator/\allowbreak{}errors.yr}).

\subsubsection*{What it means}

A \tokennolst{data} record generates its field constructor implicitly, and this declaration writes it out again.

\subsubsection*{Example}

\textit{test\_\allowbreak{}resources/\allowbreak{}default\_\allowbreak{}ctor/\allowbreak{}test10.yr}:

%% check: skip
\begin{lstlisting}[style=coloredverbatimError]
// A `@data` record asks for the ctor by its attribute, the request is implicit
// and is not redeclarable.
@data
record Foo {
    let x: i32;

    pub self;
}

fn main () {
    let a = Foo (1);
    a.x;
}
\end{lstlisting}

\begin{lstlisting}[style=bashVerb, breaklines=true, escapechar=~]
Error[E3007] : the ctor built from the fields is implicit in a data record
 --> test_resources/default_ctor/test10.yr:(7,9)
 7  ~\lstbox{\textSFxi{}}~     pub self;
    ~\lstbox{\textSFv{}}~         ^^^^
    ~\lstbox{\textSFxi{}}~ Note : declared here
    ~\lstbox{\textSFxi{}}~  --> test_resources/default_ctor/test10.yr:(4,8)
    ~\lstbox{\textSFxi{}}~  4  ~\lstbox{\textSFxi{}}~ record Foo {
    ~\lstbox{\textSFxi{}}~     ~\lstbox{\textSFv{}}~        ^^^
    ~\lstbox{\textSFxi{}}~     ~\lstbox{\textSFii{}}~~\lstbox{\textSFx{}}~~\lstbox{\textSFx{}}~~\lstbox{\textSFx{}}~~\lstbox{\textSFx{}}~~\lstbox{\textSFx{}}~~\lstbox{\textSFx{}}~~\lstbox{\textSFx{}}~
    ~\lstbox{\textSFii{}}~~\lstbox{\textSFx{}}~~\lstbox{\textSFx{}}~~\lstbox{\textSFx{}}~~\lstbox{\textSFx{}}~~\lstbox{\textSFx{}}~~\lstbox{\textSFvii{}}~~\lstbox{\textSFx{}}~
\end{lstlisting}

\subsubsection*{How to fix it}

Remove the explicit constructor: a \tokennolst{data} record always has it.

\errorcode{E3010}{a class may only define one destructor}
\label{sec:codes:e3010}

Raised during \textbf{declaration}, from \texttt{DeclareErrorMessage::\allowbreak{}MULTIPLE\_\allowbreak{}DESTRUCTOR} (\textit{src/\allowbreak{}ymirc/\allowbreak{}semantic/\allowbreak{}declarator/\allowbreak{}errors.yr}).

\subsubsection*{What it means}

A class declares more than one destructor. A destructor takes no parameter, so there is nothing to overload on.

\subsubsection*{Example}

\textit{test\_\allowbreak{}resources/\allowbreak{}diagnostics/\allowbreak{}test6.yr}:

%% check: skip
\begin{lstlisting}[style=coloredverbatimError]
class Foo {
    pub self () {}
    __dtor (mut self) {}
    __dtor (mut self) {}
}

fn main () {}
\end{lstlisting}

\begin{lstlisting}[style=bashVerb, breaklines=true, escapechar=~]
Error[E3010] : a class may only define one destructor
 --> test_resources/diagnostics/test6.yr:(4,5)
 4  ~\lstbox{\textSFxi{}}~     __dtor (mut self) {}
    ~\lstbox{\textSFv{}}~     ^^^^^^
    ~\lstbox{\textSFxi{}}~ Note :
    ~\lstbox{\textSFxi{}}~  --> test_resources/diagnostics/test6.yr:(3,5)
    ~\lstbox{\textSFxi{}}~  3  ~\lstbox{\textSFxi{}}~     __dtor (mut self) {}
    ~\lstbox{\textSFxi{}}~     ~\lstbox{\textSFv{}}~     ^^^^^^
    ~\lstbox{\textSFxi{}}~     ~\lstbox{\textSFii{}}~~\lstbox{\textSFx{}}~~\lstbox{\textSFx{}}~~\lstbox{\textSFx{}}~~\lstbox{\textSFx{}}~~\lstbox{\textSFx{}}~~\lstbox{\textSFx{}}~~\lstbox{\textSFx{}}~
    ~\lstbox{\textSFii{}}~~\lstbox{\textSFx{}}~~\lstbox{\textSFx{}}~~\lstbox{\textSFx{}}~~\lstbox{\textSFx{}}~~\lstbox{\textSFx{}}~~\lstbox{\textSFvii{}}~~\lstbox{\textSFx{}}~
\end{lstlisting}

\subsubsection*{How to fix it}

Keep one destructor and merge the bodies.

\errorcode{E3015}{a record cannot be both overlaid and packed}
\label{sec:codes:e3015}

Raised during \textbf{declaration}, from \texttt{DeclareErrorMessage::\allowbreak{}OVERLAID\_\allowbreak{}AND\_\allowbreak{}PACKED} (\textit{src/\allowbreak{}ymirc/\allowbreak{}semantic/\allowbreak{}declarator/\allowbreak{}errors.yr}).

\subsubsection*{What it means}

A record is declared both overlaid and packed. Overlaying places every field at the same offset; packing removes the padding between successive fields. There is no layout satisfying both.

\subsubsection*{Example}

\textit{test\_\allowbreak{}resources/\allowbreak{}diagnostics/\allowbreak{}test8.yr}:

%% check: skip
\begin{lstlisting}[style=coloredverbatimError]
@overlaid @packed
record Foo {
    let x: i32;
    let y: i32;
}

fn main () { let _ = Foo (1, 2); }
\end{lstlisting}

\begin{lstlisting}[style=bashVerb, breaklines=true, escapechar=~]
Error[E3015] : a record cannot be both overlaid and packed
 --> test_resources/diagnostics/test8.yr:(1,12)
 1  ~\lstbox{\textSFxi{}}~ @overlaid @packed
    ~\lstbox{\textSFv{}}~            ^^^^^^
    ~\lstbox{\textSFxi{}}~ Note :
    ~\lstbox{\textSFxi{}}~  --> test_resources/diagnostics/test8.yr:(1,2)
    ~\lstbox{\textSFxi{}}~  1  ~\lstbox{\textSFxi{}}~ @overlaid @packed
    ~\lstbox{\textSFxi{}}~     ~\lstbox{\textSFv{}}~  ^^^^^^^^
    ~\lstbox{\textSFxi{}}~     ~\lstbox{\textSFii{}}~~\lstbox{\textSFx{}}~~\lstbox{\textSFx{}}~~\lstbox{\textSFx{}}~~\lstbox{\textSFx{}}~~\lstbox{\textSFx{}}~~\lstbox{\textSFx{}}~~\lstbox{\textSFx{}}~
    ~\lstbox{\textSFxi{}}~ Note :
    ~\lstbox{\textSFxi{}}~  --> test_resources/diagnostics/test8.yr:(2,1)
    ~\lstbox{\textSFxi{}}~  2  ~\lstbox{\textSFxi{}}~ record Foo {
    ~\lstbox{\textSFxi{}}~     ~\lstbox{\textSFv{}}~ ^^^^^^
    ~\lstbox{\textSFxi{}}~     ~\lstbox{\textSFii{}}~~\lstbox{\textSFx{}}~~\lstbox{\textSFx{}}~~\lstbox{\textSFx{}}~~\lstbox{\textSFx{}}~~\lstbox{\textSFx{}}~~\lstbox{\textSFx{}}~~\lstbox{\textSFx{}}~
    ~\lstbox{\textSFii{}}~~\lstbox{\textSFx{}}~~\lstbox{\textSFx{}}~~\lstbox{\textSFx{}}~~\lstbox{\textSFx{}}~~\lstbox{\textSFx{}}~~\lstbox{\textSFvii{}}~~\lstbox{\textSFx{}}~
\end{lstlisting}

\subsubsection*{How to fix it}

Keep whichever describes the layout wanted. A union-like type is overlaid; a type mirroring a wire format is packed.

\errorcode{E3030}{unexpected declaration inside an impl block}
\label{sec:codes:e3030}

Raised during \textbf{declaration}, from \texttt{DeclareErrorMessage::\allowbreak{}UNEXPECTED\_\allowbreak{}IN\_\allowbreak{}IMPL} (\textit{src/\allowbreak{}ymirc/\allowbreak{}semantic/\allowbreak{}declarator/\allowbreak{}errors.yr}).

\subsubsection*{What it means}

A declaration appeared inside an \tokennolst{impl} block that such a block does not accept. An \tokennolst{impl} block provides the methods of a trait for a type.

\subsubsection*{Example}

\textit{test\_\allowbreak{}resources/\allowbreak{}diagnostics/\allowbreak{}test96.yr}:

%% check: skip
\begin{lstlisting}[style=coloredverbatimError]
// a template method inside an impl block
trait Shape {
    pub fn area (self)-> i32;
}

class Square {
    pub self () {}

    impl Shape {
        pub over area {T} (self)-> i32 { 1 }
    }
}

fn main () {
    let _ = copy Square ();
}
\end{lstlisting}

\begin{lstlisting}[style=bashVerb, breaklines=true, escapechar=~]
Error[E3030] : unexpected declaration inside an impl block
 --> test_resources/diagnostics/test96.yr:(10,18)
10  ~\lstbox{\textSFxi{}}~         pub over area {T} (self)-> i32 { 1 }
    ~\lstbox{\textSFv{}}~                  ^^^^
    ~\lstbox{\textSFii{}}~~\lstbox{\textSFx{}}~~\lstbox{\textSFx{}}~~\lstbox{\textSFx{}}~~\lstbox{\textSFx{}}~~\lstbox{\textSFx{}}~~\lstbox{\textSFx{}}~~\lstbox{\textSFx{}}~
\end{lstlisting}

\subsubsection*{How to fix it}

Move the declaration out of the block. A method that is not part of the trait belongs in the type body rather than in the \tokennolst{impl}.

\errorcode{E3033}{unexpected declaration inside a trait}
\label{sec:codes:e3033}

Raised during \textbf{declaration}, from \texttt{DeclareErrorMessage::\allowbreak{}UNEXPECTED\_\allowbreak{}IN\_\allowbreak{}TRAIT} (\textit{src/\allowbreak{}ymirc/\allowbreak{}semantic/\allowbreak{}declarator/\allowbreak{}errors.yr}).

\subsubsection*{What it means}

A declaration appeared inside a trait that a trait body does not accept. A trait declares methods, with or without a default body.

\subsubsection*{Example}

\textit{test\_\allowbreak{}resources/\allowbreak{}diagnostics/\allowbreak{}test97.yr}:

%% check: skip
\begin{lstlisting}[style=coloredverbatimError]
// a use declaration inside a trait
trait Shape {
    use std::io;
}

fn main () {}
\end{lstlisting}

\begin{lstlisting}[style=bashVerb, breaklines=true, escapechar=~]
Error[E3033] : unexpected declaration inside a trait
 --> test_resources/diagnostics/test97.yr:(3,5)
 3  ~\lstbox{\textSFxi{}}~     use std::io;
    ~\lstbox{\textSFv{}}~     ^^^
    ~\lstbox{\textSFii{}}~~\lstbox{\textSFx{}}~~\lstbox{\textSFx{}}~~\lstbox{\textSFx{}}~~\lstbox{\textSFx{}}~~\lstbox{\textSFx{}}~~\lstbox{\textSFx{}}~~\lstbox{\textSFx{}}~
\end{lstlisting}

\subsubsection*{How to fix it}

Move the declaration out of the trait. A trait cannot hold fields: the state a method needs belongs to the type implementing it.

\errorcode{E4011}{cannot cast derived type \errhole{} to base class type \errhole{} in that context}
\label{sec:codes:e4011}

Raised during \textbf{validation}, from \texttt{ValidateErrorMessage::\allowbreak{}CANNOT\_\allowbreak{}CASTTO\_\allowbreak{}BASE\_\allowbreak{}CLASS} (\textit{src/\allowbreak{}ymirc/\allowbreak{}semantic/\allowbreak{}validator/\allowbreak{}errors.yr}).

\subsubsection*{What it means}

A value of a derived class type is cast to one of its base classes in a context that does not allow it. The cast itself is meaningful, but not where it was written.

\subsubsection*{Example}

\textit{test\_\allowbreak{}resources/\allowbreak{}diagnostics/\allowbreak{}test129.yr}:

%% check: skip
\begin{lstlisting}[style=coloredverbatimError]
// a derived instance bound by reference to its base class in a pattern
class Base {
    pub self () {}
}

class Derived over Base {
    pub self () with super () {}
}

fn main () {
    let d = copy Derived ();
    match d {
        ref b : &Base => { let _ = b; }
        _ => {}
    }
}
\end{lstlisting}

\begin{lstlisting}[style=bashVerb, breaklines=true, escapechar=~]
Error[E4011] : cannot cast derived type &(test129::Derived) to base class type &(test129::Base) in that context
 --> test_resources/diagnostics/test129.yr:(12,11)
12  ~\lstbox{\textSFxi{}}~     match d {
    ~\lstbox{\textSFv{}}~           ^
13  ~\lstbox{\textSFxi{}}~         ref b : &Base => { let _ = b; }
    ~\lstbox{\textSFv{}}~         --- -
    ~\lstbox{\textSFxi{}}~
    = when validating pattern matching with value d of type &(test129::Derived) -- (test_resources/diagnostics/test129.yr:(12,5))
    = when validating test129::main -- (test_resources/diagnostics/test129.yr:(10,4))
    ~\lstbox{\textSFii{}}~~\lstbox{\textSFx{}}~~\lstbox{\textSFx{}}~~\lstbox{\textSFx{}}~~\lstbox{\textSFx{}}~~\lstbox{\textSFx{}}~~\lstbox{\textSFx{}}~~\lstbox{\textSFx{}}~
\end{lstlisting}

\subsubsection*{How to fix it}

Bind the value to a variable of the base class type, which performs the conversion where it is allowed, and use that variable.

\errorcode{E4017}{field \errhole{} has not yet been validated}
\label{sec:codes:e4017}

Raised during \textbf{validation}, from \texttt{ValidateErrorMessage::\allowbreak{}CLASS\_\allowbreak{}FIELD\_\allowbreak{}NOT\_\allowbreak{}VALIDATED\_\allowbreak{}YET} (\textit{src/\allowbreak{}ymirc/\allowbreak{}semantic/\allowbreak{}validator/\allowbreak{}errors.yr}).

\subsubsection*{What it means}

A field is used while the type it belongs to is still being validated, so its type is not known yet. This happens when a type refers to itself through a path the compiler has to resolve before it can finish the type.

\subsubsection*{Example}

\textit{test\_\allowbreak{}resources/\allowbreak{}lit\_\allowbreak{}class/\allowbreak{}ctors/\allowbreak{}test23.yr}:

%% check: skip
\begin{lstlisting}[style=coloredverbatimError]
class A {
    let x : i32;
    let y : i32;

    pub self ()
        with y = self.x
        , x = 9
    {}
}
\end{lstlisting}

\begin{lstlisting}[style=bashVerb, breaklines=true, escapechar=~]
Error[E4017] : field x has not yet been validated
 --> test_resources/lit_class/ctors/test23.yr:(6,23)
 6  ~\lstbox{\textSFxi{}}~         with y = self.x
    ~\lstbox{\textSFv{}}~                       ^
    ~\lstbox{\textSFxi{}}~ Note : declared here
    ~\lstbox{\textSFxi{}}~  --> test_resources/lit_class/ctors/test23.yr:(2,9)
    ~\lstbox{\textSFxi{}}~  2  ~\lstbox{\textSFxi{}}~     let x : i32;
    ~\lstbox{\textSFxi{}}~     ~\lstbox{\textSFv{}}~         ^
    ~\lstbox{\textSFxi{}}~     ~\lstbox{\textSFii{}}~~\lstbox{\textSFx{}}~~\lstbox{\textSFx{}}~~\lstbox{\textSFx{}}~~\lstbox{\textSFx{}}~~\lstbox{\textSFx{}}~~\lstbox{\textSFx{}}~~\lstbox{\textSFx{}}~
    ~\lstbox{\textSFii{}}~~\lstbox{\textSFx{}}~~\lstbox{\textSFx{}}~~\lstbox{\textSFx{}}~~\lstbox{\textSFx{}}~~\lstbox{\textSFx{}}~~\lstbox{\textSFvii{}}~~\lstbox{\textSFx{}}~
\end{lstlisting}

\subsubsection*{How to fix it}

Break the cycle: hold the self-reference behind a class (which is a reference) or a pointer, rather than by value.

\errorcode{E4018}{class type \errhole{} does not implement trait \errhole{}}
\label{sec:codes:e4018}

Raised during \textbf{validation}, from \texttt{ValidateErrorMessage::\allowbreak{}CLASS\_\allowbreak{}NOT\_\allowbreak{}IMPL} (\textit{src/\allowbreak{}ymirc/\allowbreak{}semantic/\allowbreak{}validator/\allowbreak{}errors.yr}).

\subsubsection*{What it means}

A class is used where the trait named is required, and the class does not implement it.

\subsubsection*{Example}

\textit{test\_\allowbreak{}resources/\allowbreak{}templates/\allowbreak{}explicit/\allowbreak{}test13.yr}:

%% check: skip
\begin{lstlisting}[style=coloredverbatimError]
trait R {I} {}

class B  {
    pub self () {}
    impl R!i32;
}

fn foo {T impl X!{I}, X, I of f32} () {
}

fn main () {
    foo!(&B) ();
}
\end{lstlisting}

\begin{lstlisting}[style=bashVerb, breaklines=true, escapechar=~]
Error[E4018] : class type &(test13::B) does not implement trait X!{I}
 --> test_resources/templates/explicit/test13.yr:(8,9)
 8  ~\lstbox{\textSFxi{}}~ fn foo {T impl X!{I}, X, I of f32} () {
    ~\lstbox{\textSFv{}}~         ^
    ~\lstbox{\textSFxi{}}~ Error[E4095] : incompatible types f32 and i32
    ~\lstbox{\textSFxi{}}~  --> test_resources/templates/explicit/test13.yr:(5,12)
    ~\lstbox{\textSFxi{}}~  5  ~\lstbox{\textSFxi{}}~     impl R!i32;
    ~\lstbox{\textSFxi{}}~     ~\lstbox{\textSFv{}}~            ^^^
    ~\lstbox{\textSFxi{}}~     ~\lstbox{\textSFxi{}}~ Note :
    ~\lstbox{\textSFxi{}}~     ~\lstbox{\textSFxi{}}~  --> test_resources/templates/explicit/test13.yr:(8,31)
    ~\lstbox{\textSFxi{}}~     ~\lstbox{\textSFxi{}}~  8  ~\lstbox{\textSFxi{}}~ fn foo {T impl X!{I}, X, I of f32} () {
    ~\lstbox{\textSFxi{}}~     ~\lstbox{\textSFxi{}}~     ~\lstbox{\textSFv{}}~                               ^^^
    ~\lstbox{\textSFxi{}}~     ~\lstbox{\textSFxi{}}~     ~\lstbox{\textSFii{}}~~\lstbox{\textSFx{}}~~\lstbox{\textSFx{}}~~\lstbox{\textSFx{}}~~\lstbox{\textSFx{}}~~\lstbox{\textSFx{}}~~\lstbox{\textSFx{}}~~\lstbox{\textSFx{}}~
    ~\lstbox{\textSFxi{}}~     ~\lstbox{\textSFii{}}~~\lstbox{\textSFx{}}~~\lstbox{\textSFx{}}~~\lstbox{\textSFx{}}~~\lstbox{\textSFx{}}~~\lstbox{\textSFx{}}~~\lstbox{\textSFvii{}}~~\lstbox{\textSFx{}}~
    ~\lstbox{\textSFii{}}~~\lstbox{\textSFx{}}~~\lstbox{\textSFx{}}~~\lstbox{\textSFx{}}~~\lstbox{\textSFx{}}~~\lstbox{\textSFx{}}~~\lstbox{\textSFvii{}}~~\lstbox{\textSFx{}}~
\end{lstlisting}

\subsubsection*{How to fix it}

Add an \tokennolst{impl} block for the trait to the class, or use a type that already implements it.

\errorcode{E4019}{class \errhole{} has no field named \errhole{}}
\label{sec:codes:e4019}

Raised during \textbf{validation}, from \texttt{ValidateErrorMessage::\allowbreak{}CLASS\_\allowbreak{}NO\_\allowbreak{}FIELD} (\textit{src/\allowbreak{}ymirc/\allowbreak{}semantic/\allowbreak{}validator/\allowbreak{}errors.yr}).

\subsubsection*{What it means}

A field name was used on a class that has no field of that name.

\subsubsection*{Example}

\textit{test\_\allowbreak{}resources/\allowbreak{}diagnostics/\allowbreak{}test63.yr}:

%% check: skip
\begin{lstlisting}[style=coloredverbatimError]
class Foo {
    pub self () {
        self.super ();
    }
}

fn main () { let _ = copy Foo (); }
\end{lstlisting}

\begin{lstlisting}[style=bashVerb, breaklines=true, escapechar=~]
Error[E4019] : class mut &(mut test63::Foo) has no field named super
 --> test_resources/diagnostics/test63.yr:(3,14)
 3  ~\lstbox{\textSFxi{}}~         self.super ();
    ~\lstbox{\textSFv{}}~              ^^^^^  note: class mut &(mut test63::Foo) has no method named super
    ~\lstbox{\textSFxi{}}~                  ~\lstbox{\textSFii{}}~~\lstbox{\textSFx{}}~ note: class mut &(mut test63::Foo) has no field access method named super
    ~\lstbox{\textSFxi{}}~ Error[E4173] : class &(test63::Foo) has no ancestor
    ~\lstbox{\textSFxi{}}~  --> test_resources/diagnostics/test63.yr:(3,9)
    ~\lstbox{\textSFxi{}}~  3  ~\lstbox{\textSFxi{}}~         self.super ();
    ~\lstbox{\textSFxi{}}~     ~\lstbox{\textSFv{}}~         ^^^^ ^^^^^
    ~\lstbox{\textSFxi{}}~     ~\lstbox{\textSFii{}}~~\lstbox{\textSFx{}}~~\lstbox{\textSFx{}}~~\lstbox{\textSFx{}}~~\lstbox{\textSFx{}}~~\lstbox{\textSFx{}}~~\lstbox{\textSFx{}}~~\lstbox{\textSFx{}}~
    ~\lstbox{\textSFii{}}~~\lstbox{\textSFx{}}~~\lstbox{\textSFx{}}~~\lstbox{\textSFx{}}~~\lstbox{\textSFx{}}~~\lstbox{\textSFx{}}~~\lstbox{\textSFvii{}}~~\lstbox{\textSFx{}}~
\end{lstlisting}

\subsubsection*{Notes}

When the name is a method rather than a field, the reason it cannot be reached is reported under this diagnostic, as a note carrying no code of its own:

\begin{itemize}
\item \tokennolst{method \textless{}\textgreater{} is a field method} (formerly E4061): a field method was reached through the alias operator \tokennolst{:.}.
\item \texttt{the superclass proxy of \textless{}\textgreater{} is only accessible through self} (formerly E4200): \tokennolst{super} was accessed on a value other than \tokennolst{self}.
\end{itemize}

\subsubsection*{How to fix it}

Check the spelling. A field that exists but is not visible from here gives a protection error instead, so this really means no such field. A field method of that name is reached the same way, so check that it is declared too.

\errorcode{E4020}{class \errhole{} has no method named \errhole{}}
\label{sec:codes:e4020}

Raised during \textbf{validation}, from \texttt{ValidateErrorMessage::\allowbreak{}CLASS\_\allowbreak{}NO\_\allowbreak{}METHOD} (\textit{src/\allowbreak{}ymirc/\allowbreak{}semantic/\allowbreak{}validator/\allowbreak{}errors.yr}).

\subsubsection*{What it means}

A method name was used on a class that has no method of that name.

\subsubsection*{Example}

\textit{test\_\allowbreak{}resources/\allowbreak{}lit\_\allowbreak{}class/\allowbreak{}operators/\allowbreak{}test13.yr}:

%% check: skip
\begin{lstlisting}[style=coloredverbatimError]
class A {
    pub self () {}

    pub fn opIndex (self, _ : i32)-> i32 {
        1
    }

}


class B  {
    pub self () {}
}



fn main () {
    let a = copy A ();
    let _ = a ['i'];

    let b = copy B ();
    let _ = b [1];
}
\end{lstlisting}

\begin{lstlisting}[style=bashVerb, breaklines=true, escapechar=~]
Error[E4020] : class &(test13::B) has no method named opIndex
 --> test_resources/lit_class/operators/test13.yr:(22,15)
22  ~\lstbox{\textSFxi{}}~     let _ = b [1];
    ~\lstbox{\textSFv{}}~               ^
    ~\lstbox{\textSFxi{}}~
    = when validating test13::main -- (test_resources/lit_class/operators/test13.yr:(17,4))
    ~\lstbox{\textSFii{}}~~\lstbox{\textSFx{}}~~\lstbox{\textSFx{}}~~\lstbox{\textSFx{}}~~\lstbox{\textSFx{}}~~\lstbox{\textSFx{}}~~\lstbox{\textSFx{}}~~\lstbox{\textSFx{}}~
\end{lstlisting}

\subsubsection*{How to fix it}

Check the spelling and the class. When the method comes from a trait, the class needs an \tokennolst{impl} block for it, cf. \hyperref[sec:codes:e4018]{E4018}.

\errorcode{E4024}{colliding method definitions \errhole{} and \errhole{}}
\label{sec:codes:e4024}

Raised during \textbf{validation}, from \texttt{ValidateErrorMessage::\allowbreak{}COLLIDING\_\allowbreak{}METHOD\_\allowbreak{}DEFINITION} (\textit{src/\allowbreak{}ymirc/\allowbreak{}semantic/\allowbreak{}validator/\allowbreak{}errors.yr}).

\subsubsection*{What it means}

Two methods of a class have the same name and indistinguishable signatures, cf. \hyperref[sec:codes:e4023]{E4023}.

\subsubsection*{Example}

\textit{test\_\allowbreak{}resources/\allowbreak{}default\_\allowbreak{}ctor/\allowbreak{}test6.yr}:

%% check: skip
\begin{lstlisting}[style=coloredverbatimError]
// A `self;` declares the ctor built from the fields, so a ctor written by hand
// with the same signature declares it twice.
class Foo {
   pub let x: i32;

   pub self (x: i32) with x = x + 1 {}
   pub self;
}

fn main () {
    let a = copy Foo (1);
    a.x;
}
\end{lstlisting}

\begin{lstlisting}[style=bashVerb, breaklines=true, escapechar=~]
Error[E4024] : colliding method definitions test6::Foo::self and test6::Foo::self
 --> test_resources/default_ctor/test6.yr:(6,8)
 6  ~\lstbox{\textSFxi{}}~    pub self (x: i32) with x = x + 1 {}
    ~\lstbox{\textSFv{}}~        ^^^^
 7  ~\lstbox{\textSFxi{}}~    pub self;
    ~\lstbox{\textSFv{}}~        ----
    ~\lstbox{\textSFxi{}}~
    = when validating -- (test_resources/default_ctor/test6.yr:(3,7))
    ~\lstbox{\textSFii{}}~~\lstbox{\textSFx{}}~~\lstbox{\textSFx{}}~~\lstbox{\textSFx{}}~~\lstbox{\textSFx{}}~~\lstbox{\textSFx{}}~~\lstbox{\textSFx{}}~~\lstbox{\textSFx{}}~
Error[E4104] : symbol test6::Foo has errors
 --> test_resources/default_ctor/test6.yr:(11,18)
11  ~\lstbox{\textSFxi{}}~     let a = copy Foo (1);
    ~\lstbox{\textSFv{}}~                  ^^^
 --> test_resources/default_ctor/test6.yr:(6,8)
 6  ~\lstbox{\textSFxi{}}~    pub self (x: i32) with x = x + 1 {}
    ~\lstbox{\textSFv{}}~        ----  error: colliding method definitions test6::Foo::self and test6::Foo::self
    ~\lstbox{\textSFxi{}}~
    = when validating -- (test_resources/default_ctor/test6.yr:(3,7))
    = when validating test6::main -- (test_resources/default_ctor/test6.yr:(10,4))
    ~\lstbox{\textSFii{}}~~\lstbox{\textSFx{}}~~\lstbox{\textSFx{}}~~\lstbox{\textSFx{}}~~\lstbox{\textSFx{}}~~\lstbox{\textSFx{}}~~\lstbox{\textSFx{}}~~\lstbox{\textSFx{}}~
Error[E4224] : undefined operator . for types error and x
 --> test_resources/default_ctor/test6.yr:(12,6)
12  ~\lstbox{\textSFxi{}}~     a.x;
    ~\lstbox{\textSFv{}}~      ^
    ~\lstbox{\textSFxi{}}~
    = when validating test6::main -- (test_resources/default_ctor/test6.yr:(10,4))
    ~\lstbox{\textSFii{}}~~\lstbox{\textSFx{}}~~\lstbox{\textSFx{}}~~\lstbox{\textSFx{}}~~\lstbox{\textSFx{}}~~\lstbox{\textSFx{}}~~\lstbox{\textSFx{}}~~\lstbox{\textSFx{}}~
\end{lstlisting}

\subsubsection*{How to fix it}

Rename one of them, or change the parameters.

\errorcode{E4026}{cannot construct an instance of class \errhole{} that is abstract}
\label{sec:codes:e4026}

Raised during \textbf{validation}, from \texttt{ValidateErrorMessage::\allowbreak{}CONSTRUCT\_\allowbreak{}ABSTRACT\_\allowbreak{}CLASS} (\textit{src/\allowbreak{}ymirc/\allowbreak{}semantic/\allowbreak{}validator/\allowbreak{}errors.yr}).

\subsubsection*{What it means}

An instance of an abstract class was constructed. An abstract class has methods with no body, so an instance of it would have nothing to call.

\subsubsection*{Example}

\textit{test\_\allowbreak{}resources/\allowbreak{}diagnostics/\allowbreak{}test134.yr}:

%% check: skip
\begin{lstlisting}[style=coloredverbatimError]
// an abstract class constructed
@abstract
class Shape {
    pub self () {}
}

fn main () {
    let _ = copy Shape ();
}
\end{lstlisting}

\begin{lstlisting}[style=bashVerb, breaklines=true, escapechar=~]
Error[E4026] : cannot construct an instance of class test134::Shape that is abstract
 --> test_resources/diagnostics/test134.yr:(8,18)
 8  ~\lstbox{\textSFxi{}}~     let _ = copy Shape ();
    ~\lstbox{\textSFv{}}~                  ^^^^^
    ~\lstbox{\textSFxi{}}~
    = when validating test134::main -- (test_resources/diagnostics/test134.yr:(7,4))
    ~\lstbox{\textSFii{}}~~\lstbox{\textSFx{}}~~\lstbox{\textSFx{}}~~\lstbox{\textSFx{}}~~\lstbox{\textSFx{}}~~\lstbox{\textSFx{}}~~\lstbox{\textSFx{}}~~\lstbox{\textSFx{}}~
\end{lstlisting}

\subsubsection*{How to fix it}

Construct a derived class that implements the empty methods, or give the class bodies for them and drop \tokennolst{abstract}.

\errorcode{E4027}{method \errhole{} is defined as constant}
\label{sec:codes:e4027}

Raised during \textbf{validation}, from \texttt{ValidateErrorMessage::\allowbreak{}CONST\_\allowbreak{}METHOD} (\textit{src/\allowbreak{}ymirc/\allowbreak{}semantic/\allowbreak{}validator/\allowbreak{}errors.yr}).

\subsubsection*{What it means}

A method declared constant is used where a mutable one is needed -- typically to write through the instance it is called on.

\subsubsection*{Example}

\textit{test\_\allowbreak{}resources/\allowbreak{}lit\_\allowbreak{}class/\allowbreak{}operators/\allowbreak{}test15.yr}:

%% check: skip
\begin{lstlisting}[style=coloredverbatimError]
class A {

    pub self () {}

    pub fn opIndex (self, _ : i32) {}

}

fn main () {
    let dmut a = copy A ();
    let _ = a:[0];
    let _ = a:.opIndex (0);
}
\end{lstlisting}

\begin{lstlisting}[style=bashVerb, breaklines=true, escapechar=~]
Error[E4027] : method test15::A::opIndex (_: i32)-> void is defined as constant
 --> test_resources/lit_class/operators/test15.yr:(5,12)
 5  ~\lstbox{\textSFxi{}}~     pub fn opIndex (self, _ : i32) {}
    ~\lstbox{\textSFv{}}~            ^^^^^^^
 --> test_resources/lit_class/operators/test15.yr:(11,14)
11  ~\lstbox{\textSFxi{}}~     let _ = a:[0];
    ~\lstbox{\textSFv{}}~              --  note: aliasing the value of type mut &(mut test15::A) to create a constant borrowing is prohibited
    ~\lstbox{\textSFxi{}}~
    = when validating test15::main -- (test_resources/lit_class/operators/test15.yr:(9,4))
    ~\lstbox{\textSFii{}}~~\lstbox{\textSFx{}}~~\lstbox{\textSFx{}}~~\lstbox{\textSFx{}}~~\lstbox{\textSFx{}}~~\lstbox{\textSFx{}}~~\lstbox{\textSFx{}}~~\lstbox{\textSFx{}}~
\end{lstlisting}

\subsubsection*{How to fix it}

Make the method mutable (\tokennolst{self} declared \tokennolst{mut}), or call a mutable method instead.

\errorcode{E4035}{infinite constructor redirection calls when calling \errhole{}}
\label{sec:codes:e4035}

Raised during \textbf{validation}, from \texttt{ValidateErrorMessage::\allowbreak{}CTOR\_\allowbreak{}INFINITE\_\allowbreak{}REDIRECTION} (\textit{src/\allowbreak{}ymirc/\allowbreak{}semantic/\allowbreak{}validator/\allowbreak{}errors.yr}).

\subsubsection*{What it means}

A chain of constructor redirections comes back to a constructor already in the chain, so construction would never terminate.

\subsubsection*{Example}

\textit{test\_\allowbreak{}resources/\allowbreak{}templates/\allowbreak{}class/\allowbreak{}test3.yr}:

%% check: skip
\begin{lstlisting}[style=coloredverbatimError]
class A {
    pub self (a : i32)
        with self (a)
    {}
}
\end{lstlisting}

\begin{lstlisting}[style=bashVerb, breaklines=true, escapechar=~]
Error[E4035] : infinite constructor redirection calls when calling test3::A::self (a: i32)-> mut &(mut test3::A)
 --> test_resources/templates/class/test3.yr:(3,14)
 3  ~\lstbox{\textSFxi{}}~         with self (a)
    ~\lstbox{\textSFv{}}~              ^^^^
    ~\lstbox{\textSFii{}}~~\lstbox{\textSFx{}}~~\lstbox{\textSFx{}}~~\lstbox{\textSFx{}}~~\lstbox{\textSFx{}}~~\lstbox{\textSFx{}}~~\lstbox{\textSFx{}}~~\lstbox{\textSFx{}}~
\end{lstlisting}

\subsubsection*{How to fix it}

Break the cycle: one constructor of the chain has to initialize the fields itself instead of redirecting.

\errorcode{E4044}{field must have a referenceable name}
\label{sec:codes:e4044}

Raised during \textbf{validation}, from \texttt{ValidateErrorMessage::\allowbreak{}CLASS\_\allowbreak{}FIELD\_\allowbreak{}NO\_\allowbreak{}NAME} (\textit{src/\allowbreak{}ymirc/\allowbreak{}semantic/\allowbreak{}validator/\allowbreak{}errors.yr}).

\subsubsection*{What it means}

A field was declared without a name that can be referred to.

\subsubsection*{Example}

\textit{test\_\allowbreak{}resources/\allowbreak{}diagnostics/\allowbreak{}test100.yr}:

%% check: skip
\begin{lstlisting}[style=coloredverbatimError]
// a class field named _
class Point {
    let _ : i32 = 1;
    pub self () {}
}

fn main () {
    let _ = copy Point ();
}
\end{lstlisting}

\begin{lstlisting}[style=bashVerb, breaklines=true, escapechar=~]
Error[E4044] : field must have a referenceable name
 --> test_resources/diagnostics/test100.yr:(3,9)
 3  ~\lstbox{\textSFxi{}}~     let _ : i32 = 1;
    ~\lstbox{\textSFv{}}~         ^         -
    ~\lstbox{\textSFxi{}}~
    = when validating -- (test_resources/diagnostics/test100.yr:(2,7))
    ~\lstbox{\textSFii{}}~~\lstbox{\textSFx{}}~~\lstbox{\textSFx{}}~~\lstbox{\textSFx{}}~~\lstbox{\textSFx{}}~~\lstbox{\textSFx{}}~~\lstbox{\textSFx{}}~~\lstbox{\textSFx{}}~
Error[E4104] : symbol test100::Point has errors
 --> test_resources/diagnostics/test100.yr:(8,18)
 8  ~\lstbox{\textSFxi{}}~     let _ = copy Point ();
    ~\lstbox{\textSFv{}}~                  ^^^^^
 --> test_resources/diagnostics/test100.yr:(3,9)
 3  ~\lstbox{\textSFxi{}}~     let _ : i32 = 1;
    ~\lstbox{\textSFv{}}~         -         -  error: field must have a referenceable name
    ~\lstbox{\textSFxi{}}~
    = when validating -- (test_resources/diagnostics/test100.yr:(2,7))
    = when validating test100::main -- (test_resources/diagnostics/test100.yr:(7,4))
    ~\lstbox{\textSFii{}}~~\lstbox{\textSFx{}}~~\lstbox{\textSFx{}}~~\lstbox{\textSFx{}}~~\lstbox{\textSFx{}}~~\lstbox{\textSFx{}}~~\lstbox{\textSFx{}}~~\lstbox{\textSFx{}}~
\end{lstlisting}

\subsubsection*{How to fix it}

Give the field a name.

\errorcode{E4051}{cannot perform a direct call of the method \errhole{} that is empty}
\label{sec:codes:e4051}

Raised during \textbf{validation}, from \texttt{ValidateErrorMessage::\allowbreak{}EMPTY\_\allowbreak{}METHOD\_\allowbreak{}CALL} (\textit{src/\allowbreak{}ymirc/\allowbreak{}semantic/\allowbreak{}validator/\allowbreak{}errors.yr}).

\subsubsection*{What it means}

A method with no body was called directly rather than through the virtual dispatch of an instance. An empty method has no code to jump to.

\subsubsection*{Example}

\textit{test\_\allowbreak{}resources/\allowbreak{}diagnostics/\allowbreak{}test30.yr}:

%% check: skip
\begin{lstlisting}[style=coloredverbatimError]
@abstract
class Foo {
    pub self () {}
    pub fn bar (self);
}

class Bar over Foo {
    pub self () {}
    pub over bar (self) { self.super.bar (); }
}

fn main () { let _ = copy Bar (); }
\end{lstlisting}

\begin{lstlisting}[style=bashVerb, breaklines=true, escapechar=~]
Note : when using uniform function call syntax
 --> test_resources/diagnostics/test30.yr:(9,38)
 9  ~\lstbox{\textSFxi{}}~     pub over bar (self) { self.super.bar (); }
    ~\lstbox{\textSFv{}}~              ---                     ^^^
    ~\lstbox{\textSFxi{}}~                ~\lstbox{\textSFii{}}~~\lstbox{\textSFx{}}~ note: defined here
    ~\lstbox{\textSFxi{}}~                                        ~\lstbox{\textSFii{}}~~\lstbox{\textSFx{}}~ error: test30::Bar::bar is a method and must be called from an instance
    ~\lstbox{\textSFxi{}}~
    = when validating test30::Bar::bar -- (test_resources/diagnostics/test30.yr:(9,14))
    = when validating -- (test_resources/diagnostics/test30.yr:(7,7))
    ~\lstbox{\textSFii{}}~~\lstbox{\textSFx{}}~~\lstbox{\textSFx{}}~~\lstbox{\textSFx{}}~~\lstbox{\textSFx{}}~~\lstbox{\textSFx{}}~~\lstbox{\textSFx{}}~~\lstbox{\textSFx{}}~
Error[E4019] : class &(test30::Foo) has no field named bar
 --> test_resources/diagnostics/test30.yr:(9,38)
 9  ~\lstbox{\textSFxi{}}~     pub over bar (self) { self.super.bar (); }
    ~\lstbox{\textSFv{}}~                                ----- ^^^
    ~\lstbox{\textSFxi{}}~                                    ~\lstbox{\textSFii{}}~~\lstbox{\textSFx{}}~ note: from an proxied instance of class &(test30::Foo) that is abstract
    ~\lstbox{\textSFxi{}}~                                        ~\lstbox{\textSFii{}}~~\lstbox{\textSFx{}}~ note: class &(test30::Foo) has no field access method named bar
 --> test_resources/diagnostics/test30.yr:(4,12)
 4  ~\lstbox{\textSFxi{}}~     pub fn bar (self);
    ~\lstbox{\textSFv{}}~            ---  error: cannot perform a direct call of the method test30::Foo::bar ()-> void that is empty
    ~\lstbox{\textSFxi{}}~
    = when validating test30::Bar::bar -- (test_resources/diagnostics/test30.yr:(9,14))
    = when validating -- (test_resources/diagnostics/test30.yr:(7,7))
    ~\lstbox{\textSFii{}}~~\lstbox{\textSFx{}}~~\lstbox{\textSFx{}}~~\lstbox{\textSFx{}}~~\lstbox{\textSFx{}}~~\lstbox{\textSFx{}}~~\lstbox{\textSFx{}}~~\lstbox{\textSFx{}}~
\end{lstlisting}

\subsubsection*{How to fix it}

Call it on an instance whose dynamic type provides a body, or give the method a body.

\errorcode{E4060}{field access method must return a value}
\label{sec:codes:e4060}

Raised during \textbf{validation}, from \texttt{ValidateErrorMessage::\allowbreak{}FIELD\_\allowbreak{}METHOD\_\allowbreak{}ACCESS\_\allowbreak{}NO\_\allowbreak{}RETURN} (\textit{src/\allowbreak{}ymirc/\allowbreak{}semantic/\allowbreak{}validator/\allowbreak{}errors.yr}).

\subsubsection*{What it means}

A field access method has no return value. A field access method stands in for reading a field, so it has to produce one.

\subsubsection*{Example}

\textit{test\_\allowbreak{}resources/\allowbreak{}diagnostics/\allowbreak{}test34.yr}:

%% check: skip
\begin{lstlisting}[style=coloredverbatimError]
class Foo {
    pub self () {}
    @field pub fn bar (self) {}
}

fn main () { let f = copy Foo (); f.bar; }
\end{lstlisting}

\begin{lstlisting}[style=bashVerb, breaklines=true, escapechar=~]
Error[E4060] : field access method must return a value
 --> test_resources/diagnostics/test34.yr:(3,16)
 3  ~\lstbox{\textSFxi{}}~     @field pub fn bar (self) {}
    ~\lstbox{\textSFv{}}~      -----     ^^
    ~\lstbox{\textSFxi{}}~
    = when validating test34::Foo::bar -- (test_resources/diagnostics/test34.yr:(3,19))
    = when validating -- (test_resources/diagnostics/test34.yr:(1,7))
    ~\lstbox{\textSFii{}}~~\lstbox{\textSFx{}}~~\lstbox{\textSFx{}}~~\lstbox{\textSFx{}}~~\lstbox{\textSFx{}}~~\lstbox{\textSFx{}}~~\lstbox{\textSFx{}}~~\lstbox{\textSFx{}}~
Error[E4104] : symbol test34::Foo has errors
 --> test_resources/diagnostics/test34.yr:(6,27)
 6  ~\lstbox{\textSFxi{}}~ fn main () { let f = copy Foo (); f.bar; }
    ~\lstbox{\textSFv{}}~                           ^^^
 --> test_resources/diagnostics/test34.yr:(3,6)
 3  ~\lstbox{\textSFxi{}}~     @field pub fn bar (self) {}
    ~\lstbox{\textSFv{}}~      -----     --  error: field access method must return a value
    ~\lstbox{\textSFxi{}}~
    = when validating test34::Foo::bar -- (test_resources/diagnostics/test34.yr:(3,19))
    = when validating -- (test_resources/diagnostics/test34.yr:(1,7))
    = when validating test34::main -- (test_resources/diagnostics/test34.yr:(6,4))
    ~\lstbox{\textSFii{}}~~\lstbox{\textSFx{}}~~\lstbox{\textSFx{}}~~\lstbox{\textSFx{}}~~\lstbox{\textSFx{}}~~\lstbox{\textSFx{}}~~\lstbox{\textSFx{}}~~\lstbox{\textSFx{}}~
Error[E4224] : undefined operator . for types error and bar
 --> test_resources/diagnostics/test34.yr:(6,36)
 6  ~\lstbox{\textSFxi{}}~ fn main () { let f = copy Foo (); f.bar; }
    ~\lstbox{\textSFv{}}~                                    ^
    ~\lstbox{\textSFxi{}}~
    = when validating test34::main -- (test_resources/diagnostics/test34.yr:(6,4))
    ~\lstbox{\textSFii{}}~~\lstbox{\textSFx{}}~~\lstbox{\textSFx{}}~~\lstbox{\textSFx{}}~~\lstbox{\textSFx{}}~~\lstbox{\textSFx{}}~~\lstbox{\textSFx{}}~~\lstbox{\textSFx{}}~
\end{lstlisting}

\subsubsection*{How to fix it}

Return the value the pseudo-field stands for.

\errorcode{E4062}{field assign method must be mutable}
\label{sec:codes:e4062}

Raised during \textbf{validation}, from \texttt{ValidateErrorMessage::\allowbreak{}FIELD\_\allowbreak{}METHOD\_\allowbreak{}ASSIGN\_\allowbreak{}CONST} (\textit{src/\allowbreak{}ymirc/\allowbreak{}semantic/\allowbreak{}validator/\allowbreak{}errors.yr}).

\subsubsection*{What it means}

A field assign method is declared constant. It stands in for writing a field, which modifies the instance.

\subsubsection*{Example}

\textit{test\_\allowbreak{}resources/\allowbreak{}diagnostics/\allowbreak{}test35.yr}:

%% check: skip
\begin{lstlisting}[style=coloredverbatimError]
class Foo {
    let mut x: i32 = 0;
    pub self () {}
    @field pub fn bar (self, v: i32) { self.x = v; }
}

fn main () { let dmut f = copy Foo (); f:.bar = 1; }
\end{lstlisting}

\begin{lstlisting}[style=bashVerb, breaklines=true, escapechar=~]
Error[E4062] : field assign method must be mutable
 --> test_resources/diagnostics/test35.yr:(4,16)
 4  ~\lstbox{\textSFxi{}}~     @field pub fn bar (self, v: i32) { self.x = v; }
    ~\lstbox{\textSFv{}}~      -----     ^^
    ~\lstbox{\textSFxi{}}~
    = when validating test35::Foo::bar -- (test_resources/diagnostics/test35.yr:(4,19))
    = when validating -- (test_resources/diagnostics/test35.yr:(1,7))
    ~\lstbox{\textSFii{}}~~\lstbox{\textSFx{}}~~\lstbox{\textSFx{}}~~\lstbox{\textSFx{}}~~\lstbox{\textSFx{}}~~\lstbox{\textSFx{}}~~\lstbox{\textSFx{}}~~\lstbox{\textSFx{}}~
Error[E4104] : symbol test35::Foo has errors
 --> test_resources/diagnostics/test35.yr:(7,32)
 7  ~\lstbox{\textSFxi{}}~ fn main () { let dmut f = copy Foo (); f:.bar = 1; }
    ~\lstbox{\textSFv{}}~                                ^^^
 --> test_resources/diagnostics/test35.yr:(4,6)
 4  ~\lstbox{\textSFxi{}}~     @field pub fn bar (self, v: i32) { self.x = v; }
    ~\lstbox{\textSFv{}}~      -----     --  error: field assign method must be mutable
    ~\lstbox{\textSFxi{}}~
    = when validating test35::Foo::bar -- (test_resources/diagnostics/test35.yr:(4,19))
    = when validating -- (test_resources/diagnostics/test35.yr:(1,7))
    = when validating test35::main -- (test_resources/diagnostics/test35.yr:(7,4))
    ~\lstbox{\textSFii{}}~~\lstbox{\textSFx{}}~~\lstbox{\textSFx{}}~~\lstbox{\textSFx{}}~~\lstbox{\textSFx{}}~~\lstbox{\textSFx{}}~~\lstbox{\textSFx{}}~~\lstbox{\textSFx{}}~
Error[E4015] : alias operator :. is only usable on class, record, entity, map or generator types, not f
 --> test_resources/diagnostics/test35.yr:(7,40)
 7  ~\lstbox{\textSFxi{}}~ fn main () { let dmut f = copy Foo (); f:.bar = 1; }
    ~\lstbox{\textSFv{}}~                                        ^
    ~\lstbox{\textSFxi{}}~
    = when validating test35::main -- (test_resources/diagnostics/test35.yr:(7,4))
    ~\lstbox{\textSFii{}}~~\lstbox{\textSFx{}}~~\lstbox{\textSFx{}}~~\lstbox{\textSFx{}}~~\lstbox{\textSFx{}}~~\lstbox{\textSFx{}}~~\lstbox{\textSFx{}}~~\lstbox{\textSFx{}}~
\end{lstlisting}

\subsubsection*{How to fix it}

Declare \tokennolst{self} mutable on the method.

\errorcode{E4063}{field assign method must not return a value}
\label{sec:codes:e4063}

Raised during \textbf{validation}, from \texttt{ValidateErrorMessage::\allowbreak{}FIELD\_\allowbreak{}METHOD\_\allowbreak{}ASSIGN\_\allowbreak{}WITH\_\allowbreak{}RETURN} (\textit{src/\allowbreak{}ymirc/\allowbreak{}semantic/\allowbreak{}validator/\allowbreak{}errors.yr}).

\subsubsection*{What it means}

A field assign method returns a value. It stands in for an assignment, which produces nothing.

\subsubsection*{Example}

\textit{test\_\allowbreak{}resources/\allowbreak{}diagnostics/\allowbreak{}test36.yr}:

%% check: skip
\begin{lstlisting}[style=coloredverbatimError]
class Foo {
    let mut x: i32 = 0;
    pub self () {}
    @field pub fn bar (mut self, v: i32)-> i32 { self.x = v; v }
}

fn main () { let dmut f = copy Foo (); f:.bar = 1; }
\end{lstlisting}

\begin{lstlisting}[style=bashVerb, breaklines=true, escapechar=~]
Error[E4063] : field assign method must not return a value
 --> test_resources/diagnostics/test36.yr:(4,16)
 4  ~\lstbox{\textSFxi{}}~     @field pub fn bar (mut self, v: i32)-> i32 { self.x = v; v }
    ~\lstbox{\textSFv{}}~      -----     ^^
    ~\lstbox{\textSFxi{}}~
    = when validating test36::Foo::bar -- (test_resources/diagnostics/test36.yr:(4,19))
    = when validating -- (test_resources/diagnostics/test36.yr:(1,7))
    ~\lstbox{\textSFii{}}~~\lstbox{\textSFx{}}~~\lstbox{\textSFx{}}~~\lstbox{\textSFx{}}~~\lstbox{\textSFx{}}~~\lstbox{\textSFx{}}~~\lstbox{\textSFx{}}~~\lstbox{\textSFx{}}~
Error[E4104] : symbol test36::Foo has errors
 --> test_resources/diagnostics/test36.yr:(7,32)
 7  ~\lstbox{\textSFxi{}}~ fn main () { let dmut f = copy Foo (); f:.bar = 1; }
    ~\lstbox{\textSFv{}}~                                ^^^
 --> test_resources/diagnostics/test36.yr:(4,6)
 4  ~\lstbox{\textSFxi{}}~     @field pub fn bar (mut self, v: i32)-> i32 { self.x = v; v }
    ~\lstbox{\textSFv{}}~      -----     --  error: field assign method must not return a value
    ~\lstbox{\textSFxi{}}~
    = when validating test36::Foo::bar -- (test_resources/diagnostics/test36.yr:(4,19))
    = when validating -- (test_resources/diagnostics/test36.yr:(1,7))
    = when validating test36::main -- (test_resources/diagnostics/test36.yr:(7,4))
    ~\lstbox{\textSFii{}}~~\lstbox{\textSFx{}}~~\lstbox{\textSFx{}}~~\lstbox{\textSFx{}}~~\lstbox{\textSFx{}}~~\lstbox{\textSFx{}}~~\lstbox{\textSFx{}}~~\lstbox{\textSFx{}}~
Error[E4015] : alias operator :. is only usable on class, record, entity, map or generator types, not f
 --> test_resources/diagnostics/test36.yr:(7,40)
 7  ~\lstbox{\textSFxi{}}~ fn main () { let dmut f = copy Foo (); f:.bar = 1; }
    ~\lstbox{\textSFv{}}~                                        ^
    ~\lstbox{\textSFxi{}}~
    = when validating test36::main -- (test_resources/diagnostics/test36.yr:(7,4))
    ~\lstbox{\textSFii{}}~~\lstbox{\textSFx{}}~~\lstbox{\textSFx{}}~~\lstbox{\textSFx{}}~~\lstbox{\textSFx{}}~~\lstbox{\textSFx{}}~~\lstbox{\textSFx{}}~~\lstbox{\textSFx{}}~
\end{lstlisting}

\subsubsection*{How to fix it}

Remove the return value.

\errorcode{E4065}{field method \errhole{} is mutable}
\label{sec:codes:e4065}

Raised during \textbf{validation}, from \texttt{ValidateErrorMessage::\allowbreak{}FIELD\_\allowbreak{}METHOD\_\allowbreak{}MUTABLE\_\allowbreak{}ALIAS} (\textit{src/\allowbreak{}ymirc/\allowbreak{}semantic/\allowbreak{}validator/\allowbreak{}errors.yr}).

\subsubsection*{What it means}

A field method is mutable and is being used where a constant one is required -- typically through a value borrowed as constant.

\subsubsection*{Example}

\textit{test\_\allowbreak{}resources/\allowbreak{}lit\_\allowbreak{}class/\allowbreak{}methods/\allowbreak{}test8.yr}:

%% check: skip
\begin{lstlisting}[style=coloredverbatimError]
record A {
    let mut _x : i32 = 0;

    pub self () {}

    @field
    pub fn foo (self)-> i32 {
        self._x
    }

    @field
    pub fn foo (mut self, i : i32) {
        self._x = i;
    }

}


fn main () {
    let a = A ();
    a.foo += 1;
}
\end{lstlisting}

\begin{lstlisting}[style=bashVerb, breaklines=true, escapechar=~]
Error[E4065] : field method test8::A::foo (i: i32)-> void is mutable
 --> test_resources/lit_class/methods/test8.yr:(12,12)
12  ~\lstbox{\textSFxi{}}~     pub fn foo (mut self, i : i32) {
    ~\lstbox{\textSFv{}}~            ^^^
    ~\lstbox{\textSFxi{}}~ Error[E4082] : left operand of type test8::A is immutable
    ~\lstbox{\textSFxi{}}~  --> test_resources/lit_class/methods/test8.yr:(21,5)
    ~\lstbox{\textSFxi{}}~ 21  ~\lstbox{\textSFxi{}}~     a.foo += 1;
    ~\lstbox{\textSFxi{}}~     ~\lstbox{\textSFv{}}~     ^^
    ~\lstbox{\textSFxi{}}~     ~\lstbox{\textSFii{}}~~\lstbox{\textSFx{}}~~\lstbox{\textSFx{}}~~\lstbox{\textSFx{}}~~\lstbox{\textSFx{}}~~\lstbox{\textSFx{}}~~\lstbox{\textSFx{}}~~\lstbox{\textSFx{}}~
    ~\lstbox{\textSFii{}}~~\lstbox{\textSFx{}}~~\lstbox{\textSFx{}}~~\lstbox{\textSFx{}}~~\lstbox{\textSFx{}}~~\lstbox{\textSFx{}}~~\lstbox{\textSFvii{}}~~\lstbox{\textSFx{}}~
\end{lstlisting}

\subsubsection*{How to fix it}

Use a constant field method for reading, and reserve the mutable one for writing.

\errorcode{E4066}{field method cannot have more than one parameter}
\label{sec:codes:e4066}

Raised during \textbf{validation}, from \texttt{ValidateErrorMessage::\allowbreak{}FIELD\_\allowbreak{}METHOD\_\allowbreak{}PARAMS} (\textit{src/\allowbreak{}ymirc/\allowbreak{}semantic/\allowbreak{}validator/\allowbreak{}errors.yr}).

\subsubsection*{What it means}

A field method declares more than one parameter besides \tokennolst{self}. A field access method takes none and a field assign method takes exactly the value being assigned.

\subsubsection*{Example}

\textit{test\_\allowbreak{}resources/\allowbreak{}diagnostics/\allowbreak{}test37.yr}:

%% check: skip
\begin{lstlisting}[style=coloredverbatimError]
class Foo {
    pub self () {}
    @field pub fn bar (self, a: i32, b: i32)-> i32 { a + b }
}

fn main () { let f = copy Foo (); f.bar; }
\end{lstlisting}

\begin{lstlisting}[style=bashVerb, breaklines=true, escapechar=~]
Error[E4066] : field method cannot have more than one parameter
 --> test_resources/diagnostics/test37.yr:(3,16)
 3  ~\lstbox{\textSFxi{}}~     @field pub fn bar (self, a: i32, b: i32)-> i32 { a + b }
    ~\lstbox{\textSFv{}}~      -----     ^^
    ~\lstbox{\textSFxi{}}~
    = when validating test37::Foo::bar -- (test_resources/diagnostics/test37.yr:(3,19))
    = when validating -- (test_resources/diagnostics/test37.yr:(1,7))
    ~\lstbox{\textSFii{}}~~\lstbox{\textSFx{}}~~\lstbox{\textSFx{}}~~\lstbox{\textSFx{}}~~\lstbox{\textSFx{}}~~\lstbox{\textSFx{}}~~\lstbox{\textSFx{}}~~\lstbox{\textSFx{}}~
Error[E4104] : symbol test37::Foo has errors
 --> test_resources/diagnostics/test37.yr:(6,27)
 6  ~\lstbox{\textSFxi{}}~ fn main () { let f = copy Foo (); f.bar; }
    ~\lstbox{\textSFv{}}~                           ^^^
 --> test_resources/diagnostics/test37.yr:(3,6)
 3  ~\lstbox{\textSFxi{}}~     @field pub fn bar (self, a: i32, b: i32)-> i32 { a + b }
    ~\lstbox{\textSFv{}}~      -----     --  error: field method cannot have more than one parameter
    ~\lstbox{\textSFxi{}}~
    = when validating test37::Foo::bar -- (test_resources/diagnostics/test37.yr:(3,19))
    = when validating -- (test_resources/diagnostics/test37.yr:(1,7))
    = when validating test37::main -- (test_resources/diagnostics/test37.yr:(6,4))
    ~\lstbox{\textSFii{}}~~\lstbox{\textSFx{}}~~\lstbox{\textSFx{}}~~\lstbox{\textSFx{}}~~\lstbox{\textSFx{}}~~\lstbox{\textSFx{}}~~\lstbox{\textSFx{}}~~\lstbox{\textSFx{}}~
Error[E4224] : undefined operator . for types error and bar
 --> test_resources/diagnostics/test37.yr:(6,36)
 6  ~\lstbox{\textSFxi{}}~ fn main () { let f = copy Foo (); f.bar; }
    ~\lstbox{\textSFv{}}~                                    ^
    ~\lstbox{\textSFxi{}}~
    = when validating test37::main -- (test_resources/diagnostics/test37.yr:(6,4))
    ~\lstbox{\textSFii{}}~~\lstbox{\textSFx{}}~~\lstbox{\textSFx{}}~~\lstbox{\textSFx{}}~~\lstbox{\textSFx{}}~~\lstbox{\textSFx{}}~~\lstbox{\textSFx{}}~~\lstbox{\textSFx{}}~
\end{lstlisting}

\subsubsection*{How to fix it}

Reduce the parameter list. A method genuinely needing several arguments is an ordinary method, not a field one.

\errorcode{E4079}{the field \errhole{} is not public, and has no initial value to be constructed from}
\label{sec:codes:e4079}

Raised during \textbf{validation}, from \texttt{ValidateErrorMessage::\allowbreak{}GEN\_\allowbreak{}CTOR\_\allowbreak{}HIDDEN\_\allowbreak{}FIELD\_\allowbreak{}NO\_\allowbreak{}VALUE} (\textit{src/\allowbreak{}ymirc/\allowbreak{}semantic/\allowbreak{}validator/\allowbreak{}errors.yr}).

\subsubsection*{What it means}

The constructor built from the fields cannot initialize the reported field: it is not public, so a caller could not pass it, and it has no initial value to fall back on.

\subsubsection*{Example}

\textit{test\_\allowbreak{}resources/\allowbreak{}data\_\allowbreak{}record/\allowbreak{}ctors/\allowbreak{}test3.yr}:

%% check: skip
\begin{lstlisting}[style=coloredverbatimError]
@data
record Foo {
    prv let _x: i32;
    let y: i32;
}

fn main() {
    let a = Foo(1);
    a.y;
}
\end{lstlisting}

\begin{lstlisting}[style=bashVerb, breaklines=true, escapechar=~]
Error[E4079] : the field _x is not public, and has no initial value to be constructed from
 --> test_resources/data_record/ctors/test3.yr:(3,13)
 3  ~\lstbox{\textSFxi{}}~     prv let _x: i32;
    ~\lstbox{\textSFv{}}~             ^^
    ~\lstbox{\textSFxi{}}~
    = when validating -- (test_resources/data_record/ctors/test3.yr:(2,8))
    ~\lstbox{\textSFii{}}~~\lstbox{\textSFx{}}~~\lstbox{\textSFx{}}~~\lstbox{\textSFx{}}~~\lstbox{\textSFx{}}~~\lstbox{\textSFx{}}~~\lstbox{\textSFx{}}~~\lstbox{\textSFx{}}~
Error[E4104] : symbol test3::Foo has errors
 --> test_resources/data_record/ctors/test3.yr:(8,13)
 8  ~\lstbox{\textSFxi{}}~     let a = Foo(1);
    ~\lstbox{\textSFv{}}~             ^^^
 --> test_resources/data_record/ctors/test3.yr:(3,13)
 3  ~\lstbox{\textSFxi{}}~     prv let _x: i32;
    ~\lstbox{\textSFv{}}~             --  error: the field _x is not public, and has no initial value to be constructed from
    ~\lstbox{\textSFxi{}}~
    = when validating -- (test_resources/data_record/ctors/test3.yr:(2,8))
    = when validating test3::main -- (test_resources/data_record/ctors/test3.yr:(7,4))
    ~\lstbox{\textSFii{}}~~\lstbox{\textSFx{}}~~\lstbox{\textSFx{}}~~\lstbox{\textSFx{}}~~\lstbox{\textSFx{}}~~\lstbox{\textSFx{}}~~\lstbox{\textSFx{}}~~\lstbox{\textSFx{}}~
Error[E4224] : undefined operator . for types error and y
 --> test_resources/data_record/ctors/test3.yr:(9,6)
 9  ~\lstbox{\textSFxi{}}~     a.y;
    ~\lstbox{\textSFv{}}~      ^
    ~\lstbox{\textSFxi{}}~
    = when validating test3::main -- (test_resources/data_record/ctors/test3.yr:(7,4))
    ~\lstbox{\textSFii{}}~~\lstbox{\textSFx{}}~~\lstbox{\textSFx{}}~~\lstbox{\textSFx{}}~~\lstbox{\textSFx{}}~~\lstbox{\textSFx{}}~~\lstbox{\textSFx{}}~~\lstbox{\textSFx{}}~
\end{lstlisting}

\subsubsection*{How to fix it}

Give the field an initial value, make it public, or write the constructor explicitly.

\errorcode{E4090}{method \errhole{} implicitly overrides method \errhole{}}
\label{sec:codes:e4090}

Raised during \textbf{validation}, from \texttt{ValidateErrorMessage::\allowbreak{}IMPLICIT\_\allowbreak{}OVERRIDE} (\textit{src/\allowbreak{}ymirc/\allowbreak{}semantic/\allowbreak{}validator/\allowbreak{}errors.yr}).

\subsubsection*{What it means}

A method has the same signature as one of an ancestor class but is not marked \tokennolst{override}. Ymir requires the marker so that adding a method to a base class cannot silently change what a derived class does.

\subsubsection*{Example}

\textit{test\_\allowbreak{}resources/\allowbreak{}diagnostics/\allowbreak{}test42.yr}:

%% check: skip
\begin{lstlisting}[style=coloredverbatimError]
class A {
    pub self () {}
    pub fn foo (self) {}
}

class B over A {
    pub self () {}
    pub fn foo (self) {}
}

fn main () { let _ = copy B (); }
\end{lstlisting}

\begin{lstlisting}[style=bashVerb, breaklines=true, escapechar=~]
Error[E4090] : method test42::B::foo ()-> void implicitly overrides method test42::A::foo ()-> void
 --> test_resources/diagnostics/test42.yr:(3,12)
 3  ~\lstbox{\textSFxi{}}~     pub fn foo (self) {}
    ~\lstbox{\textSFv{}}~            ^^^
 --> test_resources/diagnostics/test42.yr:(8,12)
 8  ~\lstbox{\textSFxi{}}~     pub fn foo (self) {}
    ~\lstbox{\textSFv{}}~            ---
    ~\lstbox{\textSFxi{}}~
    = when validating -- (test_resources/diagnostics/test42.yr:(6,7))
    ~\lstbox{\textSFii{}}~~\lstbox{\textSFx{}}~~\lstbox{\textSFx{}}~~\lstbox{\textSFx{}}~~\lstbox{\textSFx{}}~~\lstbox{\textSFx{}}~~\lstbox{\textSFx{}}~~\lstbox{\textSFx{}}~
Error[E4104] : symbol test42::B has errors
 --> test_resources/diagnostics/test42.yr:(11,27)
11  ~\lstbox{\textSFxi{}}~ fn main () { let _ = copy B (); }
    ~\lstbox{\textSFv{}}~                           ^
 --> test_resources/diagnostics/test42.yr:(3,12)
 3  ~\lstbox{\textSFxi{}}~     pub fn foo (self) {}
    ~\lstbox{\textSFv{}}~            ---  error: method test42::B::foo ()-> void implicitly overrides method test42::A::foo ()-> void
    ~\lstbox{\textSFxi{}}~
    = when validating -- (test_resources/diagnostics/test42.yr:(6,7))
    = when validating test42::main -- (test_resources/diagnostics/test42.yr:(11,4))
    ~\lstbox{\textSFii{}}~~\lstbox{\textSFx{}}~~\lstbox{\textSFx{}}~~\lstbox{\textSFx{}}~~\lstbox{\textSFx{}}~~\lstbox{\textSFx{}}~~\lstbox{\textSFx{}}~~\lstbox{\textSFx{}}~
\end{lstlisting}

\subsubsection*{How to fix it}

Add \tokennolst{override} if overriding was intended, or rename the method if it was meant to be a new one.

\errorcode{E4091}{implicit override of method \errhole{} by trait impl}
\label{sec:codes:e4091}

Raised during \textbf{validation}, from \texttt{ValidateErrorMessage::\allowbreak{}IMPLICIT\_\allowbreak{}OVERRIDE\_\allowbreak{}BY\_\allowbreak{}TRAIT} (\textit{src/\allowbreak{}ymirc/\allowbreak{}semantic/\allowbreak{}validator/\allowbreak{}errors.yr}).

\subsubsection*{What it means}

A trait \tokennolst{impl} provides a method that overrides an ancestor method, without the override being declared.

\subsubsection*{Example}

\textit{test\_\allowbreak{}resources/\allowbreak{}diagnostics/\allowbreak{}test107.yr}:

%% check: skip
\begin{lstlisting}[style=coloredverbatimError]
// a trait implemented by a class whose ancestor already has a method of that name
trait Shape {
    pub fn area (self)-> i32 { 1 }
}

class Base {
    pub self () {}
    pub fn area (self)-> i32 { 2 }
}

class Square over Base {
    pub self () with super () {}
    impl Shape;
}

fn main () {
    let _ = copy Square ();
}
\end{lstlisting}

\begin{lstlisting}[style=bashVerb, breaklines=true, escapechar=~]
Error[E4091] : implicit override of method test107::Base::area by trait impl
 --> test_resources/diagnostics/test107.yr:(13,5)
13  ~\lstbox{\textSFxi{}}~     impl Shape;
    ~\lstbox{\textSFv{}}~     ^^^^
 --> test_resources/diagnostics/test107.yr:(3,12)
 3  ~\lstbox{\textSFxi{}}~     pub fn area (self)-> i32 { 1 }
    ~\lstbox{\textSFv{}}~            ----
    ~\lstbox{\textSFxi{}}~
    = when validating -- (test_resources/diagnostics/test107.yr:(11,7))
    ~\lstbox{\textSFii{}}~~\lstbox{\textSFx{}}~~\lstbox{\textSFx{}}~~\lstbox{\textSFx{}}~~\lstbox{\textSFx{}}~~\lstbox{\textSFx{}}~~\lstbox{\textSFx{}}~~\lstbox{\textSFx{}}~
Error[E4104] : symbol test107::Square has errors
 --> test_resources/diagnostics/test107.yr:(17,18)
17  ~\lstbox{\textSFxi{}}~     let _ = copy Square ();
    ~\lstbox{\textSFv{}}~                  ^^^^^^
 --> test_resources/diagnostics/test107.yr:(3,12)
 3  ~\lstbox{\textSFxi{}}~     pub fn area (self)-> i32 { 1 }
    ~\lstbox{\textSFv{}}~            ----
   ...
13  ~\lstbox{\textSFxi{}}~     impl Shape;
    ~\lstbox{\textSFv{}}~     ----  error: implicit override of method test107::Base::area by trait impl
    ~\lstbox{\textSFxi{}}~
    = when validating -- (test_resources/diagnostics/test107.yr:(11,7))
    = when validating test107::main -- (test_resources/diagnostics/test107.yr:(16,4))
    ~\lstbox{\textSFii{}}~~\lstbox{\textSFx{}}~~\lstbox{\textSFx{}}~~\lstbox{\textSFx{}}~~\lstbox{\textSFx{}}~~\lstbox{\textSFx{}}~~\lstbox{\textSFx{}}~~\lstbox{\textSFx{}}~
\end{lstlisting}

\subsubsection*{How to fix it}

Mark the method \tokennolst{override} in the \tokennolst{impl}, or rename it.

\errorcode{E4093}{trait \errhole{} is implemented multiple times by \errhole{}}
\label{sec:codes:e4093}

Raised during \textbf{validation}, from \texttt{ValidateErrorMessage::\allowbreak{}IMPL\_\allowbreak{}MULTIPLE\_\allowbreak{}TIMES} (\textit{src/\allowbreak{}ymirc/\allowbreak{}semantic/\allowbreak{}validator/\allowbreak{}errors.yr}).

\subsubsection*{What it means}

The same trait is implemented twice by one type, so a call to one of its methods would have two candidates.

\subsubsection*{Example}

\textit{test\_\allowbreak{}resources/\allowbreak{}diagnostics/\allowbreak{}test43.yr}:

%% check: skip
\begin{lstlisting}[style=coloredverbatimError]
trait T { pub fn f (self); }

class Foo {
    pub self () {}
    impl T { pub over f (self) {} }
    impl T { pub over f (self) {} }
}

fn main () {}
\end{lstlisting}

\begin{lstlisting}[style=bashVerb, breaklines=true, escapechar=~]
Error[E4093] : trait test43::T is implemented multiple times by test43::Foo
 --> test_resources/diagnostics/test43.yr:(6,5)
 6  ~\lstbox{\textSFxi{}}~     impl T { pub over f (self) {} }
    ~\lstbox{\textSFv{}}~     ^^^^
    ~\lstbox{\textSFxi{}}~ Note : here
    ~\lstbox{\textSFxi{}}~  --> test_resources/diagnostics/test43.yr:(1,7)
    ~\lstbox{\textSFxi{}}~  1  ~\lstbox{\textSFxi{}}~ trait T { pub fn f (self); }
    ~\lstbox{\textSFxi{}}~     ~\lstbox{\textSFv{}}~       ^
    ~\lstbox{\textSFxi{}}~     ~\lstbox{\textSFii{}}~~\lstbox{\textSFx{}}~~\lstbox{\textSFx{}}~~\lstbox{\textSFx{}}~~\lstbox{\textSFx{}}~~\lstbox{\textSFx{}}~~\lstbox{\textSFx{}}~~\lstbox{\textSFx{}}~
    ~\lstbox{\textSFii{}}~~\lstbox{\textSFx{}}~~\lstbox{\textSFx{}}~~\lstbox{\textSFx{}}~~\lstbox{\textSFx{}}~~\lstbox{\textSFx{}}~~\lstbox{\textSFvii{}}~~\lstbox{\textSFx{}}~
\end{lstlisting}

\subsubsection*{How to fix it}

Merge the two \tokennolst{impl} blocks.

\errorcode{E4094}{impl statement must be followed by a trait, not \errhole{}}
\label{sec:codes:e4094}

Raised during \textbf{validation}, from \texttt{ValidateErrorMessage::\allowbreak{}IMPL\_\allowbreak{}NO\_\allowbreak{}TRAIT} (\textit{src/\allowbreak{}ymirc/\allowbreak{}semantic/\allowbreak{}validator/\allowbreak{}errors.yr}).

\subsubsection*{What it means}

An \tokennolst{impl} statement names something that is not a trait.

\subsubsection*{Example}

\textit{test\_\allowbreak{}resources/\allowbreak{}diagnostics/\allowbreak{}test44.yr}:

%% check: skip
\begin{lstlisting}[style=coloredverbatimError]
class Bar { pub self () {} }

class Foo {
    pub self () {}
    impl Bar {}
}

fn main () {}
\end{lstlisting}

\begin{lstlisting}[style=bashVerb, breaklines=true, escapechar=~]
Error[E4094] : impl statement must be followed by a trait, not test44::Bar
 --> test_resources/diagnostics/test44.yr:(1,7)
 1  ~\lstbox{\textSFxi{}}~ class Bar { pub self () {} }
    ~\lstbox{\textSFv{}}~       ^^^
    ~\lstbox{\textSFii{}}~~\lstbox{\textSFx{}}~~\lstbox{\textSFx{}}~~\lstbox{\textSFx{}}~~\lstbox{\textSFx{}}~~\lstbox{\textSFx{}}~~\lstbox{\textSFx{}}~~\lstbox{\textSFx{}}~
\end{lstlisting}

\subsubsection*{How to fix it}

Name a trait. To add methods to a type without a trait, declare them in the type body.

\errorcode{E4101}{the base class \errhole{} is marked as final}
\label{sec:codes:e4101}

Raised during \textbf{validation}, from \texttt{ValidateErrorMessage::\allowbreak{}INHERIT\_\allowbreak{}FINAL\_\allowbreak{}CLASS} (\textit{src/\allowbreak{}ymirc/\allowbreak{}semantic/\allowbreak{}validator/\allowbreak{}errors.yr}).

\subsubsection*{What it means}

A class inherits from a class marked \tokennolst{final}. \tokennolst{final} says the class is a leaf of its hierarchy.

\subsubsection*{Example}

\textit{test\_\allowbreak{}resources/\allowbreak{}diagnostics/\allowbreak{}test45.yr}:

%% check: skip
\begin{lstlisting}[style=coloredverbatimError]
@final
class A { pub self () {} }

class B over A { pub self () {} }

fn main () { let _ = copy B (); }
\end{lstlisting}

\begin{lstlisting}[style=bashVerb, breaklines=true, escapechar=~]
Error[E4101] : the base class test45::A is marked as final
 --> test_resources/diagnostics/test45.yr:(4,14)
 4  ~\lstbox{\textSFxi{}}~ class B over A { pub self () {} }
    ~\lstbox{\textSFv{}}~              ^
    ~\lstbox{\textSFxi{}}~
    = when validating -- (test_resources/diagnostics/test45.yr:(4,7))
    ~\lstbox{\textSFii{}}~~\lstbox{\textSFx{}}~~\lstbox{\textSFx{}}~~\lstbox{\textSFx{}}~~\lstbox{\textSFx{}}~~\lstbox{\textSFx{}}~~\lstbox{\textSFx{}}~~\lstbox{\textSFx{}}~
Error[E4104] : symbol test45::B has errors
 --> test_resources/diagnostics/test45.yr:(6,27)
 4  ~\lstbox{\textSFxi{}}~ class B over A { pub self () {} }
    ~\lstbox{\textSFv{}}~              -  error: the base class test45::A is marked as final
 5  ~\lstbox{\textSFxi{}}~
 6  ~\lstbox{\textSFxi{}}~ fn main () { let _ = copy B (); }
    ~\lstbox{\textSFv{}}~                           ^
    ~\lstbox{\textSFxi{}}~
    = when validating -- (test_resources/diagnostics/test45.yr:(4,7))
    = when validating test45::main -- (test_resources/diagnostics/test45.yr:(6,4))
    ~\lstbox{\textSFii{}}~~\lstbox{\textSFx{}}~~\lstbox{\textSFx{}}~~\lstbox{\textSFx{}}~~\lstbox{\textSFx{}}~~\lstbox{\textSFx{}}~~\lstbox{\textSFx{}}~~\lstbox{\textSFx{}}~
\end{lstlisting}

\subsubsection*{How to fix it}

Remove \tokennolst{final} from the base class, or inherit from one of its ancestors.

\errorcode{E4102}{the base of a class must be a class, not a \errhole{}}
\label{sec:codes:e4102}

Raised during \textbf{validation}, from \texttt{ValidateErrorMessage::\allowbreak{}INHERIT\_\allowbreak{}NO\_\allowbreak{}CLASS} (\textit{src/\allowbreak{}ymirc/\allowbreak{}semantic/\allowbreak{}validator/\allowbreak{}errors.yr}).

\subsubsection*{What it means}

The base of a class is not a class -- a record, a trait or a plain type was named instead.

\subsubsection*{Example}

\textit{test\_\allowbreak{}resources/\allowbreak{}diagnostics/\allowbreak{}test46.yr}:

%% check: skip
\begin{lstlisting}[style=coloredverbatimError]
record R { let x: i32; }

class B over R { pub self () {} }

fn main () { let _ = copy B (); }
\end{lstlisting}

\begin{lstlisting}[style=bashVerb, breaklines=true, escapechar=~]
Error[E4102] : the base of a class must be a class, not a test46::R
 --> test_resources/diagnostics/test46.yr:(3,14)
 3  ~\lstbox{\textSFxi{}}~ class B over R { pub self () {} }
    ~\lstbox{\textSFv{}}~              ^
    ~\lstbox{\textSFxi{}}~
    = when validating -- (test_resources/diagnostics/test46.yr:(3,7))
    ~\lstbox{\textSFii{}}~~\lstbox{\textSFx{}}~~\lstbox{\textSFx{}}~~\lstbox{\textSFx{}}~~\lstbox{\textSFx{}}~~\lstbox{\textSFx{}}~~\lstbox{\textSFx{}}~~\lstbox{\textSFx{}}~
Error[E4104] : symbol test46::B has errors
 --> test_resources/diagnostics/test46.yr:(5,27)
 3  ~\lstbox{\textSFxi{}}~ class B over R { pub self () {} }
    ~\lstbox{\textSFv{}}~              -  error: the base of a class must be a class, not a test46::R
 4  ~\lstbox{\textSFxi{}}~
 5  ~\lstbox{\textSFxi{}}~ fn main () { let _ = copy B (); }
    ~\lstbox{\textSFv{}}~                           ^
    ~\lstbox{\textSFxi{}}~
    = when validating -- (test_resources/diagnostics/test46.yr:(3,7))
    = when validating test46::main -- (test_resources/diagnostics/test46.yr:(5,4))
    ~\lstbox{\textSFii{}}~~\lstbox{\textSFx{}}~~\lstbox{\textSFx{}}~~\lstbox{\textSFx{}}~~\lstbox{\textSFx{}}~~\lstbox{\textSFx{}}~~\lstbox{\textSFx{}}~~\lstbox{\textSFx{}}~
\end{lstlisting}

\subsubsection*{How to fix it}

Inherit from a class. To reuse a trait, use an \tokennolst{impl} block; to reuse a record, hold it as a field.

\errorcode{E4103}{method \errhole{} cannot be inline as it is virtual(overridable or inherited)}
\label{sec:codes:e4103}

Raised during \textbf{validation}, from \texttt{ValidateErrorMessage::\allowbreak{}INLINE\_\allowbreak{}VIRTUAL\_\allowbreak{}METHOD} (\textit{src/\allowbreak{}ymirc/\allowbreak{}semantic/\allowbreak{}validator/\allowbreak{}errors.yr}).

\subsubsection*{What it means}

A method is marked \tokennolst{inline} and is virtual. A virtual call is resolved at run time through the vtable, so there is no single body to inline.

\subsubsection*{Example}

\textit{test\_\allowbreak{}resources/\allowbreak{}diagnostics/\allowbreak{}test47.yr}:

%% check: skip
\begin{lstlisting}[style=coloredverbatimError]
class A {
    pub self () {}
    @inline pub fn foo (self) {}
}

fn main () { let _ = copy A (); }
\end{lstlisting}

\begin{lstlisting}[style=bashVerb, breaklines=true, escapechar=~]
Error[E4103] : method test47::A::foo cannot be inline as it is virtual(overridable or inherited)
 --> test_resources/diagnostics/test47.yr:(3,6)
 3  ~\lstbox{\textSFxi{}}~     @inline pub fn foo (self) {}
    ~\lstbox{\textSFv{}}~      ^^^^^^
    ~\lstbox{\textSFxi{}}~
    = when validating test47::A::foo -- (test_resources/diagnostics/test47.yr:(3,20))
    = when validating -- (test_resources/diagnostics/test47.yr:(1,7))
    ~\lstbox{\textSFii{}}~~\lstbox{\textSFx{}}~~\lstbox{\textSFx{}}~~\lstbox{\textSFx{}}~~\lstbox{\textSFx{}}~~\lstbox{\textSFx{}}~~\lstbox{\textSFx{}}~~\lstbox{\textSFx{}}~
Error[E4104] : symbol test47::A has errors
 --> test_resources/diagnostics/test47.yr:(6,27)
 6  ~\lstbox{\textSFxi{}}~ fn main () { let _ = copy A (); }
    ~\lstbox{\textSFv{}}~                           ^
 --> test_resources/diagnostics/test47.yr:(3,6)
 3  ~\lstbox{\textSFxi{}}~     @inline pub fn foo (self) {}
    ~\lstbox{\textSFv{}}~      ------  error: method test47::A::foo cannot be inline as it is virtual(overridable or inherited)
    ~\lstbox{\textSFxi{}}~
    = when validating test47::A::foo -- (test_resources/diagnostics/test47.yr:(3,20))
    = when validating -- (test_resources/diagnostics/test47.yr:(1,7))
    = when validating test47::main -- (test_resources/diagnostics/test47.yr:(6,4))
    ~\lstbox{\textSFii{}}~~\lstbox{\textSFx{}}~~\lstbox{\textSFx{}}~~\lstbox{\textSFx{}}~~\lstbox{\textSFx{}}~~\lstbox{\textSFx{}}~~\lstbox{\textSFx{}}~~\lstbox{\textSFx{}}~
\end{lstlisting}

\subsubsection*{How to fix it}

Drop \tokennolst{inline}, or make the method \tokennolst{final} so it is no longer overridable.

\errorcode{E4133}{field \errhole{} is initialized multiple times}
\label{sec:codes:e4133}

Raised during \textbf{validation}, from \texttt{ValidateErrorMessage::\allowbreak{}MULTIPLE\_\allowbreak{}FIELD\_\allowbreak{}INIT} (\textit{src/\allowbreak{}ymirc/\allowbreak{}semantic/\allowbreak{}validator/\allowbreak{}errors.yr}).

\subsubsection*{What it means}

A constructor assigns the same field twice.

\subsubsection*{Example}

\textit{test\_\allowbreak{}resources/\allowbreak{}union/\allowbreak{}test3.yr}:

%% check: skip
\begin{lstlisting}[style=coloredverbatimError]
@overlaid
record A {
    let mut a : i32;
    let mut b : i64 = 0;

    pub self ()
        with a = 12
        , a = 9
    {}
}

fn main () {
    let _ = A ();
}
\end{lstlisting}

\begin{lstlisting}[style=bashVerb, breaklines=true, escapechar=~]
Error[E4133] : field a is initialized multiple times
 --> test_resources/union/test3.yr:(8,11)
 8  ~\lstbox{\textSFxi{}}~         , a = 9
    ~\lstbox{\textSFv{}}~           ^
    ~\lstbox{\textSFxi{}}~ Note : here
    ~\lstbox{\textSFxi{}}~  --> test_resources/union/test3.yr:(7,14)
    ~\lstbox{\textSFxi{}}~  7  ~\lstbox{\textSFxi{}}~         with a = 12
    ~\lstbox{\textSFxi{}}~     ~\lstbox{\textSFv{}}~              ^
    ~\lstbox{\textSFxi{}}~     ~\lstbox{\textSFii{}}~~\lstbox{\textSFx{}}~~\lstbox{\textSFx{}}~~\lstbox{\textSFx{}}~~\lstbox{\textSFx{}}~~\lstbox{\textSFx{}}~~\lstbox{\textSFx{}}~~\lstbox{\textSFx{}}~
    ~\lstbox{\textSFii{}}~~\lstbox{\textSFx{}}~~\lstbox{\textSFx{}}~~\lstbox{\textSFx{}}~~\lstbox{\textSFx{}}~~\lstbox{\textSFx{}}~~\lstbox{\textSFvii{}}~~\lstbox{\textSFx{}}~
\end{lstlisting}

\subsubsection*{How to fix it}

Keep one assignment. When the two came from a field constructor and an explicit assignment, drop the explicit one.

\errorcode{E4141}{method \errhole{} is defined as mutable}
\label{sec:codes:e4141}

Raised during \textbf{validation}, from \texttt{ValidateErrorMessage::\allowbreak{}MUTABLE\_\allowbreak{}METHOD} (\textit{src/\allowbreak{}ymirc/\allowbreak{}semantic/\allowbreak{}validator/\allowbreak{}errors.yr}).

\subsubsection*{What it means}

A mutable method is called on a value that is not mutable.

\subsubsection*{Example}

\textit{test\_\allowbreak{}resources/\allowbreak{}diagnostics/\allowbreak{}test55.yr}:

%% check: skip
\begin{lstlisting}[style=coloredverbatimError]
class Foo {
    let mut x: i32 = 0;
    pub self () {}
    pub fn bar (mut self) { self.x = 1; }
}

fn main () {
    let f = copy Foo ();
    f.bar ();
}
\end{lstlisting}

\begin{lstlisting}[style=bashVerb, breaklines=true, escapechar=~]
Error[E4141] : method test55::Foo::bar ()-> void is defined as mutable
 --> test_resources/diagnostics/test55.yr:(4,12)
 4  ~\lstbox{\textSFxi{}}~     pub fn bar (mut self) { self.x = 1; }
    ~\lstbox{\textSFv{}}~            ^^^
    ~\lstbox{\textSFxi{}}~ Note : implicit alias of type &(test55::Foo) is prohibited, it implicitly gives up on borrowings of mutable values
    ~\lstbox{\textSFxi{}}~  --> test_resources/diagnostics/test55.yr:(9,5)
    ~\lstbox{\textSFxi{}}~  9  ~\lstbox{\textSFxi{}}~     f.bar ();
    ~\lstbox{\textSFxi{}}~     ~\lstbox{\textSFv{}}~     ^
    ~\lstbox{\textSFxi{}}~     ~\lstbox{\textSFii{}}~~\lstbox{\textSFx{}}~~\lstbox{\textSFx{}}~~\lstbox{\textSFx{}}~~\lstbox{\textSFx{}}~~\lstbox{\textSFx{}}~~\lstbox{\textSFx{}}~~\lstbox{\textSFx{}}~
    ~\lstbox{\textSFii{}}~~\lstbox{\textSFx{}}~~\lstbox{\textSFx{}}~~\lstbox{\textSFx{}}~~\lstbox{\textSFx{}}~~\lstbox{\textSFx{}}~~\lstbox{\textSFvii{}}~~\lstbox{\textSFx{}}~
\end{lstlisting}

\subsubsection*{How to fix it}

Make the value mutable, or call a constant method instead, cf. \hyperref[sec:codes:e4082]{E4082}.

\errorcode{E4144}{class \errhole{} is not abstract but has empty methods}
\label{sec:codes:e4144}

Raised during \textbf{validation}, from \texttt{ValidateErrorMessage::\allowbreak{}NON\_\allowbreak{}ABSTRACT\_\allowbreak{}NOT\_\allowbreak{}COMPLETE} (\textit{src/\allowbreak{}ymirc/\allowbreak{}semantic/\allowbreak{}validator/\allowbreak{}errors.yr}).

\subsubsection*{What it means}

A class has methods with no body but is not declared \tokennolst{abstract}, so it could be instantiated with nothing to call.

\subsubsection*{Example}

\textit{test\_\allowbreak{}resources/\allowbreak{}implement/\allowbreak{}test3.yr}:

%% check: skip
\begin{lstlisting}[style=coloredverbatimError]
trait X {
    pub fn foo (self);
}

class A {
    pub self () {}
    impl X ;
}

record B {
    pub self () {}
    impl X ;
}
\end{lstlisting}

\begin{lstlisting}[style=bashVerb, breaklines=true, escapechar=~]
Error[E4144] : class test3::A is not abstract but has empty methods
 --> test_resources/implement/test3.yr:(5,7)
 5  ~\lstbox{\textSFxi{}}~ class A {
    ~\lstbox{\textSFv{}}~       ^  note: X::foo ()-> void is empty -- (test_resources/implement/test3.yr:(2,12))
    ~\lstbox{\textSFii{}}~~\lstbox{\textSFx{}}~~\lstbox{\textSFx{}}~~\lstbox{\textSFx{}}~~\lstbox{\textSFx{}}~~\lstbox{\textSFx{}}~~\lstbox{\textSFx{}}~~\lstbox{\textSFx{}}~
\end{lstlisting}

\subsubsection*{How to fix it}

Give the methods bodies, or declare the class \tokennolst{abstract}, cf. \hyperref[sec:codes:e4026]{E4026}.

\errorcode{E4147}{\errhole{} is not an ancestor type of \errhole{}}
\label{sec:codes:e4147}

Raised during \textbf{validation}, from \texttt{ValidateErrorMessage::\allowbreak{}NOT\_\allowbreak{}ANCESTOR} (\textit{src/\allowbreak{}ymirc/\allowbreak{}semantic/\allowbreak{}validator/\allowbreak{}errors.yr}).

\subsubsection*{What it means}

A type is used where an ancestor of another type is required, and it is not one of its ancestors.

\subsubsection*{Example}

\textit{test\_\allowbreak{}resources/\allowbreak{}regressions/\allowbreak{}throws\_\allowbreak{}without\_\allowbreak{}ampersand/\allowbreak{}test3.yr}:

%% check: skip
\begin{lstlisting}[style=coloredverbatimError]
fn main ()
    throws i32
{}

fn foo () {
    throw 12;
}
\end{lstlisting}

\begin{lstlisting}[style=bashVerb, breaklines=true, escapechar=~]
Error[E4147] : mut &(mut core::exception::Exception) is not an ancestor type of i32
 --> test_resources/regressions/throws_without_ampersand/test3.yr:(6,5)
 6  ~\lstbox{\textSFxi{}}~     throw 12;
    ~\lstbox{\textSFv{}}~     ^^^^^
    ~\lstbox{\textSFii{}}~~\lstbox{\textSFx{}}~~\lstbox{\textSFx{}}~~\lstbox{\textSFx{}}~~\lstbox{\textSFx{}}~~\lstbox{\textSFx{}}~~\lstbox{\textSFx{}}~~\lstbox{\textSFx{}}~
\end{lstlisting}

\subsubsection*{How to fix it}

Name a class the other one actually derives from.

\errorcode{E4150}{\errhole{} is not a class type}
\label{sec:codes:e4150}

Raised during \textbf{validation}, from \texttt{ValidateErrorMessage::\allowbreak{}NOT\_\allowbreak{}A\_\allowbreak{}CLASS} (\textit{src/\allowbreak{}ymirc/\allowbreak{}semantic/\allowbreak{}validator/\allowbreak{}errors.yr}).

\subsubsection*{What it means}

An operation defined only on classes was applied to something else.

\subsubsection*{Example}

\textit{test\_\allowbreak{}resources/\allowbreak{}regressions/\allowbreak{}throws\_\allowbreak{}without\_\allowbreak{}ampersand/\allowbreak{}test3.yr}:

%% check: skip
\begin{lstlisting}[style=coloredverbatimError]
fn main ()
    throws i32
{}

fn foo () {
    throw 12;
}
\end{lstlisting}

\begin{lstlisting}[style=bashVerb, breaklines=true, escapechar=~]
Error[E4150] : i32 is not a class type
 --> test_resources/regressions/throws_without_ampersand/test3.yr:(2,12)
 2  ~\lstbox{\textSFxi{}}~     throws i32
    ~\lstbox{\textSFv{}}~            ^^^
    ~\lstbox{\textSFii{}}~~\lstbox{\textSFx{}}~~\lstbox{\textSFx{}}~~\lstbox{\textSFx{}}~~\lstbox{\textSFx{}}~~\lstbox{\textSFx{}}~~\lstbox{\textSFx{}}~~\lstbox{\textSFx{}}~
\end{lstlisting}

\subsubsection*{How to fix it}

Use a class value.

\errorcode{E4162}{\errhole{} is not a trait}
\label{sec:codes:e4162}

Raised during \textbf{validation}, from \texttt{ValidateErrorMessage::\allowbreak{}NOT\_\allowbreak{}A\_\allowbreak{}TRAIT} (\textit{src/\allowbreak{}ymirc/\allowbreak{}semantic/\allowbreak{}validator/\allowbreak{}errors.yr}).

\subsubsection*{What it means}

A name is used where a trait is required and does not denote one, cf. \hyperref[sec:codes:e4094]{E4094}.

\subsubsection*{How to fix it}

Name a trait.

\errorcode{E4168}{no constructor found for class \errhole{}}
\label{sec:codes:e4168}

Raised during \textbf{validation}, from \texttt{ValidateErrorMessage::\allowbreak{}NO\_\allowbreak{}CTOR\_\allowbreak{}FOUND} (\textit{src/\allowbreak{}ymirc/\allowbreak{}semantic/\allowbreak{}validator/\allowbreak{}errors.yr}).

\subsubsection*{What it means}

A class was constructed and it declares no constructor callable that way.

\subsubsection*{Example}

\textit{test\_\allowbreak{}resources/\allowbreak{}default\_\allowbreak{}ctor/\allowbreak{}test8.yr}:

%% check: skip
\begin{lstlisting}[style=coloredverbatimError]
// The generated ctor carries the protection of the `self;` asking for it.
class Foo {
   pub let x: i32;
   prv self;

   pub fn get (self)-> i32 { self.x }
}

fn main () {
    let a = copy Foo (1);
    a.get ();
}
\end{lstlisting}

\begin{lstlisting}[style=bashVerb, breaklines=true, escapechar=~]
Error[E4168] : no constructor found for class test8::Foo
 --> test_resources/default_ctor/test8.yr:(10,18)
10  ~\lstbox{\textSFxi{}}~     let a = copy Foo (1);
    ~\lstbox{\textSFv{}}~                  ^^^
 --> test_resources/default_ctor/test8.yr:(4,8)
 4  ~\lstbox{\textSFxi{}}~    prv self;
    ~\lstbox{\textSFv{}}~        ----  error: test8::Foo::self is private in this context
    ~\lstbox{\textSFxi{}}~
    = when validating test8::main -- (test_resources/default_ctor/test8.yr:(9,4))
    ~\lstbox{\textSFii{}}~~\lstbox{\textSFx{}}~~\lstbox{\textSFx{}}~~\lstbox{\textSFx{}}~~\lstbox{\textSFx{}}~~\lstbox{\textSFx{}}~~\lstbox{\textSFx{}}~~\lstbox{\textSFx{}}~
Error[E4224] : undefined operator . for types error and get
 --> test_resources/default_ctor/test8.yr:(11,6)
11  ~\lstbox{\textSFxi{}}~     a.get ();
    ~\lstbox{\textSFv{}}~      ^
    ~\lstbox{\textSFxi{}}~
    = when validating test8::main -- (test_resources/default_ctor/test8.yr:(9,4))
    ~\lstbox{\textSFii{}}~~\lstbox{\textSFx{}}~~\lstbox{\textSFx{}}~~\lstbox{\textSFx{}}~~\lstbox{\textSFx{}}~~\lstbox{\textSFx{}}~~\lstbox{\textSFx{}}~~\lstbox{\textSFx{}}~
\end{lstlisting}

\subsubsection*{How to fix it}

Declare a constructor matching the arguments given, or call one that exists; the candidates are listed under the diagnostic.

\errorcode{E4169}{no constructor named \errhole{} found for class \errhole{}}
\label{sec:codes:e4169}

Raised during \textbf{validation}, from \texttt{ValidateErrorMessage::\allowbreak{}NO\_\allowbreak{}CTOR\_\allowbreak{}FOUND\_\allowbreak{}NAME} (\textit{src/\allowbreak{}ymirc/\allowbreak{}semantic/\allowbreak{}validator/\allowbreak{}errors.yr}).

\subsubsection*{What it means}

A named constructor was called and the class has no constructor of that name.

\subsubsection*{Example}

\textit{test\_\allowbreak{}resources/\allowbreak{}diagnostics/\allowbreak{}test62.yr}:

%% check: skip
\begin{lstlisting}[style=coloredverbatimError]
class Foo {
    pub self () {}
}

fn main () { let _ = copy Foo::named (); }
\end{lstlisting}

\begin{lstlisting}[style=bashVerb, breaklines=true, escapechar=~]
Error[E4169] : no constructor named named found for class test62::Foo
 --> test_resources/diagnostics/test62.yr:(5,30)
 5  ~\lstbox{\textSFxi{}}~ fn main () { let _ = copy Foo::named (); }
    ~\lstbox{\textSFv{}}~                              ^^
    ~\lstbox{\textSFii{}}~~\lstbox{\textSFx{}}~~\lstbox{\textSFx{}}~~\lstbox{\textSFx{}}~~\lstbox{\textSFx{}}~~\lstbox{\textSFx{}}~~\lstbox{\textSFx{}}~~\lstbox{\textSFx{}}~
\end{lstlisting}

\subsubsection*{How to fix it}

Check the spelling, or use the unnamed constructor.

\errorcode{E4172}{record or entity type \errhole{} has no size due to forward references}
\label{sec:codes:e4172}

Raised during \textbf{validation}, from \texttt{ValidateErrorMessage::\allowbreak{}NO\_\allowbreak{}SIZE\_\allowbreak{}FORWARD\_\allowbreak{}REF} (\textit{src/\allowbreak{}ymirc/\allowbreak{}semantic/\allowbreak{}validator/\allowbreak{}errors.yr}).

\subsubsection*{What it means}

The size of a record or an entity cannot be computed because one of its fields refers to a type not yet defined.

\subsubsection*{Example}

\textit{test\_\allowbreak{}resources/\allowbreak{}struct/\allowbreak{}test1.yr}:

%% check: skip
\begin{lstlisting}[style=coloredverbatimError]
in test1;

record A {
    let a : A;
}

record B {
    let c : C;
}

record C {
    let b : B;
}
\end{lstlisting}

\begin{lstlisting}[style=bashVerb, breaklines=true, escapechar=~]
Error[E4172] : record or entity type test1::A has no size due to forward references
 --> test_resources/struct/test1.yr:(3,8)
 3  ~\lstbox{\textSFxi{}}~ record A {
    ~\lstbox{\textSFv{}}~        ^
    ~\lstbox{\textSFxi{}}~ Note : for the field a
    ~\lstbox{\textSFxi{}}~  --> test_resources/struct/test1.yr:(4,9)
    ~\lstbox{\textSFxi{}}~  4  ~\lstbox{\textSFxi{}}~     let a : A;
    ~\lstbox{\textSFxi{}}~     ~\lstbox{\textSFv{}}~         ^
    ~\lstbox{\textSFxi{}}~     ~\lstbox{\textSFxi{}}~ Note : within type test1::A
    ~\lstbox{\textSFxi{}}~     ~\lstbox{\textSFxi{}}~  --> test_resources/struct/test1.yr:(3,8)
    ~\lstbox{\textSFxi{}}~     ~\lstbox{\textSFxi{}}~  3  ~\lstbox{\textSFxi{}}~ record A {
    ~\lstbox{\textSFxi{}}~     ~\lstbox{\textSFxi{}}~     ~\lstbox{\textSFv{}}~        ^
    ~\lstbox{\textSFxi{}}~     ~\lstbox{\textSFxi{}}~     ~\lstbox{\textSFii{}}~~\lstbox{\textSFx{}}~~\lstbox{\textSFx{}}~~\lstbox{\textSFx{}}~~\lstbox{\textSFx{}}~~\lstbox{\textSFx{}}~~\lstbox{\textSFx{}}~~\lstbox{\textSFx{}}~
    ~\lstbox{\textSFxi{}}~     ~\lstbox{\textSFii{}}~~\lstbox{\textSFx{}}~~\lstbox{\textSFx{}}~~\lstbox{\textSFx{}}~~\lstbox{\textSFx{}}~~\lstbox{\textSFx{}}~~\lstbox{\textSFvii{}}~~\lstbox{\textSFx{}}~
    ~\lstbox{\textSFii{}}~~\lstbox{\textSFx{}}~~\lstbox{\textSFx{}}~~\lstbox{\textSFx{}}~~\lstbox{\textSFx{}}~~\lstbox{\textSFx{}}~~\lstbox{\textSFvii{}}~~\lstbox{\textSFx{}}~
\end{lstlisting}

\subsubsection*{How to fix it}

Break the cycle: hold the offending field behind a class or a pointer, cf. \hyperref[sec:codes:e4017]{E4017}.

\errorcode{E4173}{class \errhole{} has no ancestor}
\label{sec:codes:e4173}

Raised during \textbf{validation}, from \texttt{ValidateErrorMessage::\allowbreak{}NO\_\allowbreak{}SUPER\_\allowbreak{}CLASS} (\textit{src/\allowbreak{}ymirc/\allowbreak{}semantic/\allowbreak{}validator/\allowbreak{}errors.yr}).

\subsubsection*{What it means}

\tokennolst{super} was used in a class with no base class.

\subsubsection*{Example}

\textit{test\_\allowbreak{}resources/\allowbreak{}diagnostics/\allowbreak{}test63.yr}:

%% check: skip
\begin{lstlisting}[style=coloredverbatimError]
class Foo {
    pub self () {
        self.super ();
    }
}

fn main () { let _ = copy Foo (); }
\end{lstlisting}

\begin{lstlisting}[style=bashVerb, breaklines=true, escapechar=~]
Error[E4173] : class &(test63::Foo) has no ancestor
 --> test_resources/diagnostics/test63.yr:(3,9)
 3  ~\lstbox{\textSFxi{}}~         self.super ();
    ~\lstbox{\textSFv{}}~         ^^^^ ^^^^^
    ~\lstbox{\textSFii{}}~~\lstbox{\textSFx{}}~~\lstbox{\textSFx{}}~~\lstbox{\textSFx{}}~~\lstbox{\textSFx{}}~~\lstbox{\textSFx{}}~~\lstbox{\textSFx{}}~~\lstbox{\textSFx{}}~
\end{lstlisting}

\subsubsection*{How to fix it}

Give the class a base, or remove the \tokennolst{super} access.

\errorcode{E4181}{method \errhole{} must have a body to override \errhole{}}
\label{sec:codes:e4181}

Raised during \textbf{validation}, from \texttt{ValidateErrorMessage::\allowbreak{}OVERRIDE\_\allowbreak{}EMPTY} (\textit{src/\allowbreak{}ymirc/\allowbreak{}semantic/\allowbreak{}validator/\allowbreak{}errors.yr}).

\subsubsection*{What it means}

A method marked \tokennolst{override} has no body. Overriding replaces an implementation, so it has to provide one.

\subsubsection*{Example}

\textit{test\_\allowbreak{}resources/\allowbreak{}diagnostics/\allowbreak{}test64.yr}:

%% check: skip
\begin{lstlisting}[style=coloredverbatimError]
class A {
    pub self () {}
    pub fn foo (self) {}
}

class B over A {
    pub self () {}
    pub over foo (self);
}

fn main () { let _ = copy B (); }
\end{lstlisting}

\begin{lstlisting}[style=bashVerb, breaklines=true, escapechar=~]
Error[E4181] : method test64::B::foo ()-> void must have a body to override test64::A::foo ()-> void
 --> test_resources/diagnostics/test64.yr:(8,14)
 8  ~\lstbox{\textSFxi{}}~     pub over foo (self);
    ~\lstbox{\textSFv{}}~              ^^^
    ~\lstbox{\textSFxi{}}~
    = when validating -- (test_resources/diagnostics/test64.yr:(6,7))
    ~\lstbox{\textSFii{}}~~\lstbox{\textSFx{}}~~\lstbox{\textSFx{}}~~\lstbox{\textSFx{}}~~\lstbox{\textSFx{}}~~\lstbox{\textSFx{}}~~\lstbox{\textSFx{}}~~\lstbox{\textSFx{}}~
Error[E4104] : symbol test64::B has errors
 --> test_resources/diagnostics/test64.yr:(11,27)
11  ~\lstbox{\textSFxi{}}~ fn main () { let _ = copy B (); }
    ~\lstbox{\textSFv{}}~                           ^
 --> test_resources/diagnostics/test64.yr:(8,14)
 8  ~\lstbox{\textSFxi{}}~     pub over foo (self);
    ~\lstbox{\textSFv{}}~              ---  error: method test64::B::foo ()-> void must have a body to override test64::A::foo ()-> void
    ~\lstbox{\textSFxi{}}~
    = when validating -- (test_resources/diagnostics/test64.yr:(6,7))
    = when validating test64::main -- (test_resources/diagnostics/test64.yr:(11,4))
    ~\lstbox{\textSFii{}}~~\lstbox{\textSFx{}}~~\lstbox{\textSFx{}}~~\lstbox{\textSFx{}}~~\lstbox{\textSFx{}}~~\lstbox{\textSFx{}}~~\lstbox{\textSFx{}}~~\lstbox{\textSFx{}}~
\end{lstlisting}

\subsubsection*{How to fix it}

Give the method a body, or drop \tokennolst{override} and declare it in an abstract class as an empty method.

\errorcode{E4182}{method \errhole{} must be a field method to override \errhole{}}
\label{sec:codes:e4182}

Raised during \textbf{validation}, from \texttt{ValidateErrorMessage::\allowbreak{}OVERRIDE\_\allowbreak{}FIELD\_\allowbreak{}NO\_\allowbreak{}FIELD} (\textit{src/\allowbreak{}ymirc/\allowbreak{}semantic/\allowbreak{}validator/\allowbreak{}errors.yr}).

\subsubsection*{What it means}

An ordinary method tries to override a field method. The two are reached through different syntax, so one cannot stand in for the other.

\subsubsection*{Example}

\textit{test\_\allowbreak{}resources/\allowbreak{}lit\_\allowbreak{}class/\allowbreak{}methods/\allowbreak{}test3.yr}:

%% check: skip
\begin{lstlisting}[style=coloredverbatimError]
@abstract
class A {
    let _x : i32 = 0;
    let _y : i32 = 9;

    pub self () {}

    @field
    pub fn foo (self)-> i32;
}

class B over A {
    pub self () {}

    pub over foo (self)-> i32 {
        self._x
    }
}

fn main () {
    let b : &A = copy B ();
    let _z_ = b.foo;
    let _y_ = b.foo (1);
}
\end{lstlisting}

\begin{lstlisting}[style=bashVerb, breaklines=true, escapechar=~]
Error[E4182] : method test3::B::foo ()-> i32 must be a field method to override test3::A::foo ()-> i32
 --> test_resources/lit_class/methods/test3.yr:(15,14)
15  ~\lstbox{\textSFxi{}}~     pub over foo (self)-> i32 {
    ~\lstbox{\textSFv{}}~              ^^^
 --> test_resources/lit_class/methods/test3.yr:(8,6)
 8  ~\lstbox{\textSFxi{}}~     @field
    ~\lstbox{\textSFv{}}~      -----
    ~\lstbox{\textSFxi{}}~
    = when validating -- (test_resources/lit_class/methods/test3.yr:(12,7))
    ~\lstbox{\textSFii{}}~~\lstbox{\textSFx{}}~~\lstbox{\textSFx{}}~~\lstbox{\textSFx{}}~~\lstbox{\textSFx{}}~~\lstbox{\textSFx{}}~~\lstbox{\textSFx{}}~~\lstbox{\textSFx{}}~
Error[E4104] : symbol test3::B has errors
 --> test_resources/lit_class/methods/test3.yr:(21,23)
21  ~\lstbox{\textSFxi{}}~     let b : &A = copy B ();
    ~\lstbox{\textSFv{}}~                       ^
 --> test_resources/lit_class/methods/test3.yr:(8,6)
 8  ~\lstbox{\textSFxi{}}~     @field
    ~\lstbox{\textSFv{}}~      -----
   ...
15  ~\lstbox{\textSFxi{}}~     pub over foo (self)-> i32 {
    ~\lstbox{\textSFv{}}~              ---  error: method test3::B::foo ()-> i32 must be a field method to override test3::A::foo ()-> i32
    ~\lstbox{\textSFxi{}}~
    = when validating -- (test_resources/lit_class/methods/test3.yr:(12,7))
    = when validating test3::main -- (test_resources/lit_class/methods/test3.yr:(20,4))
    ~\lstbox{\textSFii{}}~~\lstbox{\textSFx{}}~~\lstbox{\textSFx{}}~~\lstbox{\textSFx{}}~~\lstbox{\textSFx{}}~~\lstbox{\textSFx{}}~~\lstbox{\textSFx{}}~~\lstbox{\textSFx{}}~
Error[E4224] : undefined operator . for types error and foo
 --> test_resources/lit_class/methods/test3.yr:(22,16)
22  ~\lstbox{\textSFxi{}}~     let _z_ = b.foo;
    ~\lstbox{\textSFv{}}~                ^
    ~\lstbox{\textSFxi{}}~
    = when validating test3::main -- (test_resources/lit_class/methods/test3.yr:(20,4))
    ~\lstbox{\textSFii{}}~~\lstbox{\textSFx{}}~~\lstbox{\textSFx{}}~~\lstbox{\textSFx{}}~~\lstbox{\textSFx{}}~~\lstbox{\textSFx{}}~~\lstbox{\textSFx{}}~~\lstbox{\textSFx{}}~
Error[E4224] : undefined operator . for types error and foo
 --> test_resources/lit_class/methods/test3.yr:(23,16)
23  ~\lstbox{\textSFxi{}}~     let _y_ = b.foo (1);
    ~\lstbox{\textSFv{}}~                ^
    ~\lstbox{\textSFxi{}}~
    = when validating test3::main -- (test_resources/lit_class/methods/test3.yr:(20,4))
    ~\lstbox{\textSFii{}}~~\lstbox{\textSFx{}}~~\lstbox{\textSFx{}}~~\lstbox{\textSFx{}}~~\lstbox{\textSFx{}}~~\lstbox{\textSFx{}}~~\lstbox{\textSFx{}}~~\lstbox{\textSFx{}}~
\end{lstlisting}

\subsubsection*{How to fix it}

Declare the overriding method as a field method too.

\errorcode{E4183}{cannot override final method \errhole{}}
\label{sec:codes:e4183}

Raised during \textbf{validation}, from \texttt{ValidateErrorMessage::\allowbreak{}OVERRIDE\_\allowbreak{}FINAL} (\textit{src/\allowbreak{}ymirc/\allowbreak{}semantic/\allowbreak{}validator/\allowbreak{}errors.yr}).

\subsubsection*{What it means}

A method marked \tokennolst{final} in an ancestor is being overridden. \tokennolst{final} says the implementation is the last word.

\subsubsection*{Example}

\textit{test\_\allowbreak{}resources/\allowbreak{}diagnostics/\allowbreak{}test65.yr}:

%% check: skip
\begin{lstlisting}[style=coloredverbatimError]
class A {
    pub self () {}
    @final pub fn foo (self) {}
}

class B over A {
    pub self () {}
    pub over foo (self) {}
}

fn main () { let _ = copy B (); }
\end{lstlisting}

\begin{lstlisting}[style=bashVerb, breaklines=true, escapechar=~]
Error[E4183] : cannot override final method test65::A::foo ()-> void
 --> test_resources/diagnostics/test65.yr:(3,19)
 3  ~\lstbox{\textSFxi{}}~     @final pub fn foo (self) {}
    ~\lstbox{\textSFv{}}~                   ^^^
 --> test_resources/diagnostics/test65.yr:(8,14)
 8  ~\lstbox{\textSFxi{}}~     pub over foo (self) {}
    ~\lstbox{\textSFv{}}~              ---
    ~\lstbox{\textSFxi{}}~
    = when validating -- (test_resources/diagnostics/test65.yr:(6,7))
    ~\lstbox{\textSFii{}}~~\lstbox{\textSFx{}}~~\lstbox{\textSFx{}}~~\lstbox{\textSFx{}}~~\lstbox{\textSFx{}}~~\lstbox{\textSFx{}}~~\lstbox{\textSFx{}}~~\lstbox{\textSFx{}}~
Error[E4104] : symbol test65::B has errors
 --> test_resources/diagnostics/test65.yr:(11,27)
11  ~\lstbox{\textSFxi{}}~ fn main () { let _ = copy B (); }
    ~\lstbox{\textSFv{}}~                           ^
 --> test_resources/diagnostics/test65.yr:(3,19)
 3  ~\lstbox{\textSFxi{}}~     @final pub fn foo (self) {}
    ~\lstbox{\textSFv{}}~                   ---  error: cannot override final method test65::A::foo ()-> void
    ~\lstbox{\textSFxi{}}~
    = when validating -- (test_resources/diagnostics/test65.yr:(6,7))
    = when validating test65::main -- (test_resources/diagnostics/test65.yr:(11,4))
    ~\lstbox{\textSFii{}}~~\lstbox{\textSFx{}}~~\lstbox{\textSFx{}}~~\lstbox{\textSFx{}}~~\lstbox{\textSFx{}}~~\lstbox{\textSFx{}}~~\lstbox{\textSFx{}}~~\lstbox{\textSFx{}}~
\end{lstlisting}

\subsubsection*{How to fix it}

Remove \tokennolst{final} from the ancestor method, or give the new method a different name.

\errorcode{E4184}{the return type of the overriding method \errhole{} is not compatible with the return type of the ancestor method \errhole{}}
\label{sec:codes:e4184}

Raised during \textbf{validation}, from \texttt{ValidateErrorMessage::\allowbreak{}OVERRIDE\_\allowbreak{}INCOMPATIBLE\_\allowbreak{}RETURN\_\allowbreak{}TYPE} (\textit{src/\allowbreak{}ymirc/\allowbreak{}semantic/\allowbreak{}validator/\allowbreak{}errors.yr}).

\subsubsection*{What it means}

An overriding method returns a type the ancestor method's return type does not accept. A caller holding the base type has to be able to use the result.

\subsubsection*{Example}

\textit{test\_\allowbreak{}resources/\allowbreak{}implement/\allowbreak{}test8.yr}:

%% check: skip
\begin{lstlisting}[style=coloredverbatimError]
trait X {T} {
    cte if template!{T}{U of i32} {
        pub fn foo (self)-> i32;
    } else {
        pub fn foo (self)-> f32;
    }
}

class Z {
    impl X!{i32} {
        pub over foo (self)-> f32 {
            1.0f
        }
    }
}

class Y {
    impl X!{f32} {
        pub over foo (self)-> i32 {
            1
        }
    }
}
\end{lstlisting}

\begin{lstlisting}[style=bashVerb, breaklines=true, escapechar=~]
Error[E4184] : the return type of the overriding method impl::foo ()-> f32 is not compatible with the return type of the ancestor method test8::X!{i32}::X::foo ()-> i32
 --> test_resources/implement/test8.yr:(11,18)
11  ~\lstbox{\textSFxi{}}~         pub over foo (self)-> f32 {
    ~\lstbox{\textSFv{}}~                  ^^^
    ~\lstbox{\textSFxi{}}~ Error[E4095] : incompatible types i32 and f32
    ~\lstbox{\textSFxi{}}~  --> test_resources/implement/test8.yr:(3,29)
    ~\lstbox{\textSFxi{}}~  3  ~\lstbox{\textSFxi{}}~         pub fn foo (self)-> i32;
    ~\lstbox{\textSFxi{}}~     ~\lstbox{\textSFv{}}~                             ^^^
    ~\lstbox{\textSFxi{}}~    ...
    ~\lstbox{\textSFxi{}}~  --> test_resources/implement/test8.yr:(11,31)
    ~\lstbox{\textSFxi{}}~ 11  ~\lstbox{\textSFxi{}}~         pub over foo (self)-> f32 {
    ~\lstbox{\textSFxi{}}~     ~\lstbox{\textSFv{}}~                               ^^^
    ~\lstbox{\textSFxi{}}~     ~\lstbox{\textSFii{}}~~\lstbox{\textSFx{}}~~\lstbox{\textSFx{}}~~\lstbox{\textSFx{}}~~\lstbox{\textSFx{}}~~\lstbox{\textSFx{}}~~\lstbox{\textSFx{}}~~\lstbox{\textSFx{}}~
    ~\lstbox{\textSFii{}}~~\lstbox{\textSFx{}}~~\lstbox{\textSFx{}}~~\lstbox{\textSFx{}}~~\lstbox{\textSFx{}}~~\lstbox{\textSFx{}}~~\lstbox{\textSFvii{}}~~\lstbox{\textSFx{}}~
\end{lstlisting}

\subsubsection*{How to fix it}

Return the ancestor's type, or a type that converts to it.

\errorcode{E4185}{the protection \errhole{} of the overriding method \errhole{} does not match the definition in the ancestor class \errhole{}}
\label{sec:codes:e4185}

Raised during \textbf{validation}, from \texttt{ValidateErrorMessage::\allowbreak{}OVERRIDE\_\allowbreak{}MISMATCH\_\allowbreak{}PROTECTION} (\textit{src/\allowbreak{}ymirc/\allowbreak{}semantic/\allowbreak{}validator/\allowbreak{}errors.yr}).

\subsubsection*{What it means}

An overriding method has a different protection from the one it overrides. A caller reaching the method through the base class would otherwise see a protection that does not hold for the actual instance.

\subsubsection*{Example}

\textit{test\_\allowbreak{}resources/\allowbreak{}diagnostics/\allowbreak{}test66.yr}:

%% check: skip
\begin{lstlisting}[style=coloredverbatimError]
class A {
    pub self () {}
    pub fn foo (self) {}
}

class B over A {
    pub self () {}
    prv over foo (self) {}
}

fn main () { let _ = copy B (); }
\end{lstlisting}

\begin{lstlisting}[style=bashVerb, breaklines=true, escapechar=~]
Error[E4185] : the protection prv of the overriding method test66::B::foo ()-> void does not match the definition in the ancestor class pub
 --> test_resources/diagnostics/test66.yr:(3,12)
 3  ~\lstbox{\textSFxi{}}~     pub fn foo (self) {}
    ~\lstbox{\textSFv{}}~            ^^^
 --> test_resources/diagnostics/test66.yr:(8,14)
 8  ~\lstbox{\textSFxi{}}~     prv over foo (self) {}
    ~\lstbox{\textSFv{}}~              ---
    ~\lstbox{\textSFxi{}}~
    = when validating -- (test_resources/diagnostics/test66.yr:(6,7))
    ~\lstbox{\textSFii{}}~~\lstbox{\textSFx{}}~~\lstbox{\textSFx{}}~~\lstbox{\textSFx{}}~~\lstbox{\textSFx{}}~~\lstbox{\textSFx{}}~~\lstbox{\textSFx{}}~~\lstbox{\textSFx{}}~
Error[E4104] : symbol test66::B has errors
 --> test_resources/diagnostics/test66.yr:(11,27)
11  ~\lstbox{\textSFxi{}}~ fn main () { let _ = copy B (); }
    ~\lstbox{\textSFv{}}~                           ^
 --> test_resources/diagnostics/test66.yr:(3,12)
 3  ~\lstbox{\textSFxi{}}~     pub fn foo (self) {}
    ~\lstbox{\textSFv{}}~            ---  error: the protection prv of the overriding method test66::B::foo ()-> void does not match the definition in the ancestor class pub
    ~\lstbox{\textSFxi{}}~
    = when validating -- (test_resources/diagnostics/test66.yr:(6,7))
    = when validating test66::main -- (test_resources/diagnostics/test66.yr:(11,4))
    ~\lstbox{\textSFii{}}~~\lstbox{\textSFx{}}~~\lstbox{\textSFx{}}~~\lstbox{\textSFx{}}~~\lstbox{\textSFx{}}~~\lstbox{\textSFx{}}~~\lstbox{\textSFx{}}~~\lstbox{\textSFx{}}~
\end{lstlisting}

\subsubsection*{How to fix it}

Give the overriding method the same protection as the ancestor one.

\errorcode{E4186}{the throwers of the overriding method \errhole{} are not compatible with the throwers of the ancestor method \errhole{}}
\label{sec:codes:e4186}

Raised during \textbf{validation}, from \texttt{ValidateErrorMessage::\allowbreak{}OVERRIDE\_\allowbreak{}MISMATCH\_\allowbreak{}THROWERS} (\textit{src/\allowbreak{}ymirc/\allowbreak{}semantic/\allowbreak{}validator/\allowbreak{}errors.yr}).

\subsubsection*{What it means}

An overriding method can throw an exception the ancestor method does not declare. A caller holding the base type would not know to handle it.

\subsubsection*{Example}

\textit{test\_\allowbreak{}resources/\allowbreak{}diagnostics/\allowbreak{}test67.yr}:

%% check: skip
\begin{lstlisting}[style=coloredverbatimError]
class A {
    pub self () {}
    pub fn foo (self) {}
}

class B over A {
    pub self () {}
    pub over foo (self) throws AssertError {
        throw copy AssertError ("x");
    }
}

fn main () { let _ = copy B (); }
\end{lstlisting}

\begin{lstlisting}[style=bashVerb, breaklines=true, escapechar=~]
Error[E4186] : the throwers of the overriding method test67::B::foo ()-> void are not compatible with the throwers of the ancestor method test67::A::foo ()-> void
 --> test_resources/diagnostics/test67.yr:(3,12)
 3  ~\lstbox{\textSFxi{}}~     pub fn foo (self) {}
    ~\lstbox{\textSFv{}}~            ^^^
 --> test_resources/diagnostics/test67.yr:(8,32)
 8  ~\lstbox{\textSFxi{}}~     pub over foo (self) throws AssertError {
    ~\lstbox{\textSFv{}}~                                -----------  error: the method might throw an exception of type core::exception::assertion::AssertError, but that is not covered by the method of the ancestor class
    ~\lstbox{\textSFxi{}}~
    = when validating -- (test_resources/diagnostics/test67.yr:(6,7))
    ~\lstbox{\textSFii{}}~~\lstbox{\textSFx{}}~~\lstbox{\textSFx{}}~~\lstbox{\textSFx{}}~~\lstbox{\textSFx{}}~~\lstbox{\textSFx{}}~~\lstbox{\textSFx{}}~~\lstbox{\textSFx{}}~
Error[E4104] : symbol test67::B has errors
 --> test_resources/diagnostics/test67.yr:(13,27)
13  ~\lstbox{\textSFxi{}}~ fn main () { let _ = copy B (); }
    ~\lstbox{\textSFv{}}~                           ^
 --> test_resources/diagnostics/test67.yr:(3,12)
 3  ~\lstbox{\textSFxi{}}~     pub fn foo (self) {}
    ~\lstbox{\textSFv{}}~            ---  error: the throwers of the overriding method test67::B::foo ()-> void are not compatible with the throwers of the ancestor method test67::A::foo ()-> void
   ...
 8  ~\lstbox{\textSFxi{}}~     pub over foo (self) throws AssertError {
    ~\lstbox{\textSFv{}}~                                -----------  error: the method might throw an exception of type core::exception::assertion::AssertError, but that is not covered by the method of the ancestor class
    ~\lstbox{\textSFxi{}}~
    = when validating -- (test_resources/diagnostics/test67.yr:(6,7))
    = when validating test67::main -- (test_resources/diagnostics/test67.yr:(13,4))
    ~\lstbox{\textSFii{}}~~\lstbox{\textSFx{}}~~\lstbox{\textSFx{}}~~\lstbox{\textSFx{}}~~\lstbox{\textSFx{}}~~\lstbox{\textSFx{}}~~\lstbox{\textSFx{}}~~\lstbox{\textSFx{}}~
\end{lstlisting}

\subsubsection*{How to fix it}

Catch the exception inside the overriding method, or declare it on the ancestor method.

\errorcode{E4187}{trait method \errhole{} was already overridden}
\label{sec:codes:e4187}

Raised during \textbf{validation}, from \texttt{ValidateErrorMessage::\allowbreak{}OVERRIDE\_\allowbreak{}MULTIPLE\_\allowbreak{}TIMES\_\allowbreak{}TRAIT} (\textit{src/\allowbreak{}ymirc/\allowbreak{}semantic/\allowbreak{}validator/\allowbreak{}errors.yr}).

\subsubsection*{What it means}

A trait method is overridden twice, so two bodies claim the same slot.

\subsubsection*{Example}

\textit{test\_\allowbreak{}resources/\allowbreak{}diagnostics/\allowbreak{}test119.yr}:

%% check: skip
\begin{lstlisting}[style=coloredverbatimError]
// a trait method overridden twice in the same impl block
trait Shape {
    pub fn area (self)-> i32;
}

class Square {
    pub self () {}

    impl Shape {
        pub over area (self)-> i32 { 1 }
        pub over area (self)-> i32 { 2 }
    }
}

fn main () {
    let _ = copy Square ();
}
\end{lstlisting}

\begin{lstlisting}[style=bashVerb, breaklines=true, escapechar=~]
Error[E4187] : trait method impl::area ()-> i32 was already overridden
 --> test_resources/diagnostics/test119.yr:(10,18)
10  ~\lstbox{\textSFxi{}}~         pub over area (self)-> i32 { 1 }
    ~\lstbox{\textSFv{}}~                  ^^^^
11  ~\lstbox{\textSFxi{}}~         pub over area (self)-> i32 { 2 }
    ~\lstbox{\textSFv{}}~                  ----
    ~\lstbox{\textSFxi{}}~
    = when validating -- (test_resources/diagnostics/test119.yr:(6,7))
    ~\lstbox{\textSFii{}}~~\lstbox{\textSFx{}}~~\lstbox{\textSFx{}}~~\lstbox{\textSFx{}}~~\lstbox{\textSFx{}}~~\lstbox{\textSFx{}}~~\lstbox{\textSFx{}}~~\lstbox{\textSFx{}}~
Error[E4104] : symbol test119::Square has errors
 --> test_resources/diagnostics/test119.yr:(16,18)
16  ~\lstbox{\textSFxi{}}~     let _ = copy Square ();
    ~\lstbox{\textSFv{}}~                  ^^^^^^
 --> test_resources/diagnostics/test119.yr:(10,18)
10  ~\lstbox{\textSFxi{}}~         pub over area (self)-> i32 { 1 }
    ~\lstbox{\textSFv{}}~                  ----  error: trait method impl::area ()-> i32 was already overridden
    ~\lstbox{\textSFxi{}}~
    = when validating -- (test_resources/diagnostics/test119.yr:(6,7))
    = when validating test119::main -- (test_resources/diagnostics/test119.yr:(15,4))
    ~\lstbox{\textSFii{}}~~\lstbox{\textSFx{}}~~\lstbox{\textSFx{}}~~\lstbox{\textSFx{}}~~\lstbox{\textSFx{}}~~\lstbox{\textSFx{}}~~\lstbox{\textSFx{}}~~\lstbox{\textSFx{}}~
\end{lstlisting}

\subsubsection*{How to fix it}

Keep one of them.

\errorcode{E4188}{cannot override a non trait method \errhole{} with \errhole{} inside impl}
\label{sec:codes:e4188}

Raised during \textbf{validation}, from \texttt{ValidateErrorMessage::\allowbreak{}OVERRIDE\_\allowbreak{}NON\_\allowbreak{}TRAIT\_\allowbreak{}INSIDE} (\textit{src/\allowbreak{}ymirc/\allowbreak{}semantic/\allowbreak{}validator/\allowbreak{}errors.yr}).

\subsubsection*{What it means}

A method inside an \tokennolst{impl} block overrides a method that is not part of the trait being implemented. An \tokennolst{impl} block provides the trait, and nothing else.

\subsubsection*{Example}

\textit{test\_\allowbreak{}resources/\allowbreak{}diagnostics/\allowbreak{}test120.yr}:

%% check: skip
\begin{lstlisting}[style=coloredverbatimError]
// an impl block overriding a method that is not part of the trait
trait Shape {
    pub fn area (self)-> i32;
}

class Square {
    pub self () {}
    pub fn side (self)-> i32 { 1 }

    impl Shape {
        pub over area (self)-> i32 { 1 }
        pub over side (self)-> i32 { 2 }
    }
}

fn main () {
    let _ = copy Square ();
}
\end{lstlisting}

\begin{lstlisting}[style=bashVerb, breaklines=true, escapechar=~]
Error[E4188] : cannot override a non trait method test120::Square::side ()-> i32 with impl::side ()-> i32 inside impl
 --> test_resources/diagnostics/test120.yr:(8,12)
 8  ~\lstbox{\textSFxi{}}~     pub fn side (self)-> i32 { 1 }
    ~\lstbox{\textSFv{}}~            ^^^^
 --> test_resources/diagnostics/test120.yr:(12,18)
12  ~\lstbox{\textSFxi{}}~         pub over side (self)-> i32 { 2 }
    ~\lstbox{\textSFv{}}~                  ----
    ~\lstbox{\textSFxi{}}~
    = when validating -- (test_resources/diagnostics/test120.yr:(6,7))
    ~\lstbox{\textSFii{}}~~\lstbox{\textSFx{}}~~\lstbox{\textSFx{}}~~\lstbox{\textSFx{}}~~\lstbox{\textSFx{}}~~\lstbox{\textSFx{}}~~\lstbox{\textSFx{}}~~\lstbox{\textSFx{}}~
Error[E4104] : symbol test120::Square has errors
 --> test_resources/diagnostics/test120.yr:(17,18)
17  ~\lstbox{\textSFxi{}}~     let _ = copy Square ();
    ~\lstbox{\textSFv{}}~                  ^^^^^^
 --> test_resources/diagnostics/test120.yr:(8,12)
 8  ~\lstbox{\textSFxi{}}~     pub fn side (self)-> i32 { 1 }
    ~\lstbox{\textSFv{}}~            ----  error: cannot override a non trait method test120::Square::side ()-> i32 with impl::side ()-> i32 inside impl
    ~\lstbox{\textSFxi{}}~
    = when validating -- (test_resources/diagnostics/test120.yr:(6,7))
    = when validating test120::main -- (test_resources/diagnostics/test120.yr:(16,4))
    ~\lstbox{\textSFii{}}~~\lstbox{\textSFx{}}~~\lstbox{\textSFx{}}~~\lstbox{\textSFx{}}~~\lstbox{\textSFx{}}~~\lstbox{\textSFx{}}~~\lstbox{\textSFx{}}~~\lstbox{\textSFx{}}~
\end{lstlisting}

\subsubsection*{How to fix it}

Move the method out of the \tokennolst{impl} block into the type body.

\errorcode{E4189}{method \errhole{} overrides nothing}
\label{sec:codes:e4189}

Raised during \textbf{validation}, from \texttt{ValidateErrorMessage::\allowbreak{}OVERRIDE\_\allowbreak{}NOTHING} (\textit{src/\allowbreak{}ymirc/\allowbreak{}semantic/\allowbreak{}validator/\allowbreak{}errors.yr}).

\subsubsection*{What it means}

A method is marked \tokennolst{override} and no ancestor method matches it, cf. \hyperref[sec:codes:retired]{E3011}.

\subsubsection*{Example}

\textit{test\_\allowbreak{}resources/\allowbreak{}diagnostics/\allowbreak{}test68.yr}:

%% check: skip
\begin{lstlisting}[style=coloredverbatimError]
class A {
    pub self () {}
    pub over foo (self) {}
}

fn main () { let _ = copy A (); }
\end{lstlisting}

\begin{lstlisting}[style=bashVerb, breaklines=true, escapechar=~]
Error[E4189] : method test68::A::foo overrides nothing
 --> test_resources/diagnostics/test68.yr:(3,14)
 3  ~\lstbox{\textSFxi{}}~     pub over foo (self) {}
    ~\lstbox{\textSFv{}}~              ^^^
    ~\lstbox{\textSFxi{}}~
    = when validating -- (test_resources/diagnostics/test68.yr:(1,7))
    ~\lstbox{\textSFii{}}~~\lstbox{\textSFx{}}~~\lstbox{\textSFx{}}~~\lstbox{\textSFx{}}~~\lstbox{\textSFx{}}~~\lstbox{\textSFx{}}~~\lstbox{\textSFx{}}~~\lstbox{\textSFx{}}~
Error[E4104] : symbol test68::A has errors
 --> test_resources/diagnostics/test68.yr:(6,27)
 6  ~\lstbox{\textSFxi{}}~ fn main () { let _ = copy A (); }
    ~\lstbox{\textSFv{}}~                           ^
 --> test_resources/diagnostics/test68.yr:(3,14)
 3  ~\lstbox{\textSFxi{}}~     pub over foo (self) {}
    ~\lstbox{\textSFv{}}~              ---  error: method test68::A::foo overrides nothing
    ~\lstbox{\textSFxi{}}~
    = when validating -- (test_resources/diagnostics/test68.yr:(1,7))
    = when validating test68::main -- (test_resources/diagnostics/test68.yr:(6,4))
    ~\lstbox{\textSFii{}}~~\lstbox{\textSFx{}}~~\lstbox{\textSFx{}}~~\lstbox{\textSFx{}}~~\lstbox{\textSFx{}}~~\lstbox{\textSFx{}}~~\lstbox{\textSFx{}}~~\lstbox{\textSFx{}}~
\end{lstlisting}

\subsubsection*{How to fix it}

Compare the signature with the ancestor method, or drop \tokennolst{override}.

\errorcode{E4190}{field method \errhole{} cannot override the non field method \errhole{}}
\label{sec:codes:e4190}

Raised during \textbf{validation}, from \texttt{ValidateErrorMessage::\allowbreak{}OVERRIDE\_\allowbreak{}NO\_\allowbreak{}FIELD\_\allowbreak{}BY\_\allowbreak{}FIELD} (\textit{src/\allowbreak{}ymirc/\allowbreak{}semantic/\allowbreak{}validator/\allowbreak{}errors.yr}).

\subsubsection*{What it means}

A field method tries to override an ordinary method, cf. \hyperref[sec:codes:e4182]{E4182}.

\subsubsection*{Example}

\textit{test\_\allowbreak{}resources/\allowbreak{}diagnostics/\allowbreak{}test121.yr}:

%% check: skip
\begin{lstlisting}[style=coloredverbatimError]
// a field method overriding a method that is not one
class Base {
    pub self () {}
    pub fn size (self)-> i32 { 1 }
}

class Derived over Base {
    pub self () with super () {}

    @field
    pub over size (self)-> i32 { 2 }
}

fn main () {
    let _ = copy Derived ();
}
\end{lstlisting}

\begin{lstlisting}[style=bashVerb, breaklines=true, escapechar=~]
Error[E4190] : field method test121::Derived::size ()-> i32 cannot override the non field method test121::Base::size ()-> i32
 --> test_resources/diagnostics/test121.yr:(10,6)
10  ~\lstbox{\textSFxi{}}~     @field
    ~\lstbox{\textSFv{}}~      ^^^^^
11  ~\lstbox{\textSFxi{}}~     pub over size (self)-> i32 { 2 }
    ~\lstbox{\textSFv{}}~              ----
    ~\lstbox{\textSFxi{}}~
    = when validating -- (test_resources/diagnostics/test121.yr:(7,7))
    ~\lstbox{\textSFii{}}~~\lstbox{\textSFx{}}~~\lstbox{\textSFx{}}~~\lstbox{\textSFx{}}~~\lstbox{\textSFx{}}~~\lstbox{\textSFx{}}~~\lstbox{\textSFx{}}~~\lstbox{\textSFx{}}~
Error[E4104] : symbol test121::Derived has errors
 --> test_resources/diagnostics/test121.yr:(15,18)
15  ~\lstbox{\textSFxi{}}~     let _ = copy Derived ();
    ~\lstbox{\textSFv{}}~                  ^^^^^^^
 --> test_resources/diagnostics/test121.yr:(10,6)
10  ~\lstbox{\textSFxi{}}~     @field
    ~\lstbox{\textSFv{}}~      -----  error: field method test121::Derived::size ()-> i32 cannot override the non field method test121::Base::size ()-> i32
    ~\lstbox{\textSFxi{}}~
    = when validating -- (test_resources/diagnostics/test121.yr:(7,7))
    = when validating test121::main -- (test_resources/diagnostics/test121.yr:(14,4))
    ~\lstbox{\textSFii{}}~~\lstbox{\textSFx{}}~~\lstbox{\textSFx{}}~~\lstbox{\textSFx{}}~~\lstbox{\textSFx{}}~~\lstbox{\textSFx{}}~~\lstbox{\textSFx{}}~~\lstbox{\textSFx{}}~
\end{lstlisting}

\subsubsection*{How to fix it}

Declare the overriding method as an ordinary method.

\errorcode{E4192}{cannot override trait method \errhole{} with \errhole{} outside impl}
\label{sec:codes:e4192}

Raised during \textbf{validation}, from \texttt{ValidateErrorMessage::\allowbreak{}OVERRIDE\_\allowbreak{}TRAIT\_\allowbreak{}OUTSIDE} (\textit{src/\allowbreak{}ymirc/\allowbreak{}semantic/\allowbreak{}validator/\allowbreak{}errors.yr}).

\subsubsection*{What it means}

A trait method is overridden outside an \tokennolst{impl} block. The trait states which methods it provides, and their bodies belong in the block implementing it.

\subsubsection*{Example}

\textit{test\_\allowbreak{}resources/\allowbreak{}implement/\allowbreak{}test4.yr}:

%% check: skip
\begin{lstlisting}[style=coloredverbatimError]
trait X {
    pub fn foo (self);
}

@abstract
class A {
    pub self () {}
    impl X ;
}

class B over A {
    pub over foo (self) {}
}
\end{lstlisting}

\begin{lstlisting}[style=bashVerb, breaklines=true, escapechar=~]
Error[E4192] : cannot override trait method X::foo ()-> void with test4::B::foo ()-> void outside impl
 --> test_resources/implement/test4.yr:(2,12)
 2  ~\lstbox{\textSFxi{}}~     pub fn foo (self);
    ~\lstbox{\textSFv{}}~            ^^^
   ...
 --> test_resources/implement/test4.yr:(12,14)
12  ~\lstbox{\textSFxi{}}~     pub over foo (self) {}
    ~\lstbox{\textSFv{}}~              ^^^
    ~\lstbox{\textSFii{}}~~\lstbox{\textSFx{}}~~\lstbox{\textSFx{}}~~\lstbox{\textSFx{}}~~\lstbox{\textSFx{}}~~\lstbox{\textSFx{}}~~\lstbox{\textSFx{}}~~\lstbox{\textSFx{}}~
\end{lstlisting}

\subsubsection*{How to fix it}

Move the method into the \tokennolst{impl} block for that trait.

\errorcode{E4196}{ancestor cycle found when validating \errhole{}}
\label{sec:codes:e4196}

Raised during \textbf{validation}, from \texttt{ValidateErrorMessage::\allowbreak{}RECURSIVE\_\allowbreak{}ANCESTOR} (\textit{src/\allowbreak{}ymirc/\allowbreak{}semantic/\allowbreak{}validator/\allowbreak{}errors.yr}).

\subsubsection*{What it means}

A class hierarchy loops back on itself: a class is, directly or transitively, its own ancestor.

\subsubsection*{Example}

\textit{test\_\allowbreak{}resources/\allowbreak{}diagnostics/\allowbreak{}test136.yr}:

%% check: skip
\begin{lstlisting}[style=coloredverbatimError]
// a class that is its own ancestor
class Node over Node {
    pub self () {}
}

fn main () {
    let _ = copy Node ();
}
\end{lstlisting}

\begin{lstlisting}[style=bashVerb, breaklines=true, escapechar=~]
Error[E4104] : symbol test136::Node has errors
 --> test_resources/diagnostics/test136.yr:(2,17)
 2  ~\lstbox{\textSFxi{}}~ class Node over Node {
    ~\lstbox{\textSFv{}}~       ----      ^^^^  error: ancestor cycle found when validating test136::Node
    ~\lstbox{\textSFxi{}}~
    = when validating -- (test_resources/diagnostics/test136.yr:(2,7))
    = when validating -- (test_resources/diagnostics/test136.yr:(2,7))
    ~\lstbox{\textSFii{}}~~\lstbox{\textSFx{}}~~\lstbox{\textSFx{}}~~\lstbox{\textSFx{}}~~\lstbox{\textSFx{}}~~\lstbox{\textSFx{}}~~\lstbox{\textSFx{}}~~\lstbox{\textSFx{}}~
Error[E4104] : symbol test136::Node has errors
 --> test_resources/diagnostics/test136.yr:(7,18)
 7  ~\lstbox{\textSFxi{}}~     let _ = copy Node ();
    ~\lstbox{\textSFv{}}~                  ^^^^
 --> test_resources/diagnostics/test136.yr:(2,7)
 2  ~\lstbox{\textSFxi{}}~ class Node over Node {
    ~\lstbox{\textSFv{}}~       ----      ----
    ~\lstbox{\textSFxi{}}~          ~\lstbox{\textSFii{}}~~\lstbox{\textSFx{}}~ error: ancestor cycle found when validating test136::Node
    ~\lstbox{\textSFxi{}}~                    ~\lstbox{\textSFii{}}~~\lstbox{\textSFx{}}~ error: symbol test136::Node has errors
    ~\lstbox{\textSFxi{}}~
    = when validating -- (test_resources/diagnostics/test136.yr:(2,7))
    = when validating -- (test_resources/diagnostics/test136.yr:(2,7))
    = when validating test136::main -- (test_resources/diagnostics/test136.yr:(6,4))
    ~\lstbox{\textSFii{}}~~\lstbox{\textSFx{}}~~\lstbox{\textSFx{}}~~\lstbox{\textSFx{}}~~\lstbox{\textSFx{}}~~\lstbox{\textSFx{}}~~\lstbox{\textSFx{}}~~\lstbox{\textSFx{}}~
\end{lstlisting}

\subsubsection*{How to fix it}

Break the cycle in the inheritance chain.

\errorcode{E4233}{no constructor is callable with the parameters \errhole{}}
\label{sec:codes:e4233}

Raised during \textbf{validation}, from \texttt{ValidateErrorMessage::\allowbreak{}UNDEFINED\_\allowbreak{}REDIRECT\_\allowbreak{}CALL\_\allowbreak{}CTOR} (\textit{src/\allowbreak{}ymirc/\allowbreak{}semantic/\allowbreak{}validator/\allowbreak{}errors.yr}).

\subsubsection*{What it means}

A constructor redirection names no constructor callable with the arguments given.

\subsubsection*{Example}

\textit{test\_\allowbreak{}resources/\allowbreak{}diagnostics/\allowbreak{}test76.yr}:

%% check: skip
\begin{lstlisting}[style=coloredverbatimError]
class Foo {
    let x: i32;
    pub self () with self (1, 2) {}
    pub self (x: i32) with x = x {}
}

fn main () { let _ = copy Foo (); }
\end{lstlisting}

\begin{lstlisting}[style=bashVerb, breaklines=true, escapechar=~]
Error[E4233] : no constructor is callable with the parameters {1, 2}
 --> test_resources/diagnostics/test76.yr:(3,22)
 3  ~\lstbox{\textSFxi{}}~     pub self () with self (1, 2) {}
    ~\lstbox{\textSFv{}}~                      ^^^^
    ~\lstbox{\textSFxi{}}~ Error[E4226] : the call operator is not defined for new and {1: i32, 2: i32}
    ~\lstbox{\textSFxi{}}~  --> test_resources/diagnostics/test76.yr:(3,22)
    ~\lstbox{\textSFxi{}}~  3  ~\lstbox{\textSFxi{}}~     pub self () with self (1, 2) {}
    ~\lstbox{\textSFxi{}}~     ~\lstbox{\textSFv{}}~                      ^^^^
    ~\lstbox{\textSFxi{}}~     ~\lstbox{\textSFxi{}}~ Note : candidate test76::Foo::self ()-> mut &(mut test76::Foo) -- (test_resources/diagnostics/test76.yr:(3,9))
    ~\lstbox{\textSFxi{}}~     ~\lstbox{\textSFxi{}}~     ~\lstbox{\textSFxi{}}~ Error[E4212] : 0 parameters are expected, but 2 are provided
    ~\lstbox{\textSFxi{}}~     ~\lstbox{\textSFxi{}}~     ~\lstbox{\textSFxi{}}~  --> test_resources/diagnostics/test76.yr:(3,22)
    ~\lstbox{\textSFxi{}}~     ~\lstbox{\textSFxi{}}~     ~\lstbox{\textSFxi{}}~  3  ~\lstbox{\textSFxi{}}~     pub self () with self (1, 2) {}
    ~\lstbox{\textSFxi{}}~     ~\lstbox{\textSFxi{}}~     ~\lstbox{\textSFxi{}}~     ~\lstbox{\textSFv{}}~                      ^^^^
    ~\lstbox{\textSFxi{}}~     ~\lstbox{\textSFxi{}}~     ~\lstbox{\textSFxi{}}~     ~\lstbox{\textSFii{}}~~\lstbox{\textSFx{}}~~\lstbox{\textSFx{}}~~\lstbox{\textSFx{}}~~\lstbox{\textSFx{}}~~\lstbox{\textSFx{}}~~\lstbox{\textSFx{}}~~\lstbox{\textSFx{}}~
    ~\lstbox{\textSFxi{}}~     ~\lstbox{\textSFxi{}}~     ~\lstbox{\textSFii{}}~~\lstbox{\textSFx{}}~~\lstbox{\textSFx{}}~~\lstbox{\textSFx{}}~~\lstbox{\textSFx{}}~~\lstbox{\textSFx{}}~~\lstbox{\textSFvii{}}~~\lstbox{\textSFx{}}~
    ~\lstbox{\textSFxi{}}~     ~\lstbox{\textSFxi{}}~ Note : candidate test76::Foo::self (x: i32)-> mut &(mut test76::Foo) -- (test_resources/diagnostics/test76.yr:(4,9))
    ~\lstbox{\textSFxi{}}~     ~\lstbox{\textSFxi{}}~     ~\lstbox{\textSFxi{}}~ Error[E4212] : 1 parameters is expected, but 2 are provided
    ~\lstbox{\textSFxi{}}~     ~\lstbox{\textSFxi{}}~     ~\lstbox{\textSFxi{}}~  --> test_resources/diagnostics/test76.yr:(3,22)
    ~\lstbox{\textSFxi{}}~     ~\lstbox{\textSFxi{}}~     ~\lstbox{\textSFxi{}}~  3  ~\lstbox{\textSFxi{}}~     pub self () with self (1, 2) {}
    ~\lstbox{\textSFxi{}}~     ~\lstbox{\textSFxi{}}~     ~\lstbox{\textSFxi{}}~     ~\lstbox{\textSFv{}}~                      ^^^^
    ~\lstbox{\textSFxi{}}~     ~\lstbox{\textSFxi{}}~     ~\lstbox{\textSFxi{}}~     ~\lstbox{\textSFii{}}~~\lstbox{\textSFx{}}~~\lstbox{\textSFx{}}~~\lstbox{\textSFx{}}~~\lstbox{\textSFx{}}~~\lstbox{\textSFx{}}~~\lstbox{\textSFx{}}~~\lstbox{\textSFx{}}~
    ~\lstbox{\textSFxi{}}~     ~\lstbox{\textSFxi{}}~     ~\lstbox{\textSFii{}}~~\lstbox{\textSFx{}}~~\lstbox{\textSFx{}}~~\lstbox{\textSFx{}}~~\lstbox{\textSFx{}}~~\lstbox{\textSFx{}}~~\lstbox{\textSFvii{}}~~\lstbox{\textSFx{}}~
    ~\lstbox{\textSFxi{}}~     ~\lstbox{\textSFii{}}~~\lstbox{\textSFx{}}~~\lstbox{\textSFx{}}~~\lstbox{\textSFx{}}~~\lstbox{\textSFx{}}~~\lstbox{\textSFx{}}~~\lstbox{\textSFvii{}}~~\lstbox{\textSFx{}}~
    ~\lstbox{\textSFii{}}~~\lstbox{\textSFx{}}~~\lstbox{\textSFx{}}~~\lstbox{\textSFx{}}~~\lstbox{\textSFx{}}~~\lstbox{\textSFx{}}~~\lstbox{\textSFvii{}}~~\lstbox{\textSFx{}}~
\end{lstlisting}

\subsubsection*{How to fix it}

Match the parameters of an existing constructor; the candidates are listed under the diagnostic.

\errorcode{E4253}{the field \errhole{} has no initial value}
\label{sec:codes:e4253}

Raised during \textbf{validation}, from \texttt{ValidateErrorMessage::\allowbreak{}UNINIT\_\allowbreak{}FIELD} (\textit{src/\allowbreak{}ymirc/\allowbreak{}semantic/\allowbreak{}validator/\allowbreak{}errors.yr}).

\subsubsection*{What it means}

A field has no initial value where one is required -- a type whose values can be created without naming every field needs every field to have a default.

\subsubsection*{Example}

\textit{test\_\allowbreak{}resources/\allowbreak{}lit\_\allowbreak{}class/\allowbreak{}ctors/\allowbreak{}test2.yr}:

%% check: skip
\begin{lstlisting}[style=coloredverbatimError]
class A {
    let i : i32;
    pub self () {}
}
\end{lstlisting}

\begin{lstlisting}[style=bashVerb, breaklines=true, escapechar=~]
Error[E4253] : the field i has no initial value
 --> test_resources/lit_class/ctors/test2.yr:(2,9)
 2  ~\lstbox{\textSFxi{}}~     let i : i32;
    ~\lstbox{\textSFv{}}~         ^
    ~\lstbox{\textSFii{}}~~\lstbox{\textSFx{}}~~\lstbox{\textSFx{}}~~\lstbox{\textSFx{}}~~\lstbox{\textSFx{}}~~\lstbox{\textSFx{}}~~\lstbox{\textSFx{}}~~\lstbox{\textSFx{}}~
\end{lstlisting}

\subsubsection*{How to fix it}

Give the field an initial value, or construct the type through a constructor that sets it.

\errorcode{E4319}{method \errhole{} must be a generator to override the generator \errhole{}}
\label{sec:codes:e4319}

Raised during \textbf{validation}, from \texttt{ValidateErrorMessage::\allowbreak{}OVERRIDE\_\allowbreak{}GENERATOR\_\allowbreak{}NO\_\allowbreak{}GENERATOR} (\textit{src/\allowbreak{}ymirc/\allowbreak{}semantic/\allowbreak{}validator/\allowbreak{}errors.yr}).

\subsubsection*{What it means}

An ordinary method tries to override a generator, declared by \tokennolst{yield}. Being a generator is part of the contract of a method, even when the overriding method returns the same handle type, cf. \hyperref[sec:codes:e4320]{E4320}.

\subsubsection*{Example}

\textit{test\_\allowbreak{}resources/\allowbreak{}generators/\allowbreak{}test85.yr}:

%% check: skip
\begin{lstlisting}[style=coloredverbatimError]
// a method overriding a generator must be a generator
yield gen ()-> i32 {
    yield 1;
}

class A {
    pub self () {}

    pub yield f (self)-> i32 {
        yield 1;
    }
}

class B over A {
    pub self () {}

    pub over f (self)-> (yield-> i32) {
        gen ()
    }
}

fn main () {
    let _ = B ();
}
\end{lstlisting}

\begin{lstlisting}[style=bashVerb, breaklines=true, escapechar=~]
Error[E4319] : method test85::B::f ()-> (yield-> i32) must be a generator to override the generator test85::A::f ()-> (yield-> i32)
 --> test_resources/generators/test85.yr:(17,14)
17  ~\lstbox{\textSFxi{}}~     pub over f (self)-> (yield-> i32) {
    ~\lstbox{\textSFv{}}~              ^
 --> test_resources/generators/test85.yr:(9,15)
 9  ~\lstbox{\textSFxi{}}~     pub yield f (self)-> i32 {
    ~\lstbox{\textSFv{}}~               -
    ~\lstbox{\textSFxi{}}~
    = when validating test85::B -- (test_resources/generators/test85.yr:(14,7))
    ~\lstbox{\textSFii{}}~~\lstbox{\textSFx{}}~~\lstbox{\textSFx{}}~~\lstbox{\textSFx{}}~~\lstbox{\textSFx{}}~~\lstbox{\textSFx{}}~~\lstbox{\textSFx{}}~~\lstbox{\textSFx{}}~
Error[E4104] : symbol test85::B has errors
 --> test_resources/generators/test85.yr:(23,13)
23  ~\lstbox{\textSFxi{}}~     let _ = B ();
    ~\lstbox{\textSFv{}}~             ^
 --> test_resources/generators/test85.yr:(9,15)
 9  ~\lstbox{\textSFxi{}}~     pub yield f (self)-> i32 {
    ~\lstbox{\textSFv{}}~               -
   ...
17  ~\lstbox{\textSFxi{}}~     pub over f (self)-> (yield-> i32) {
    ~\lstbox{\textSFv{}}~              -  error: method test85::B::f ()-> (yield-> i32) must be a generator to override the generator test85::A::f ()-> (yield-> i32)
    ~\lstbox{\textSFxi{}}~
    = when validating test85::B -- (test_resources/generators/test85.yr:(14,7))
    = when validating test85::main -- (test_resources/generators/test85.yr:(22,4))
    ~\lstbox{\textSFii{}}~~\lstbox{\textSFx{}}~~\lstbox{\textSFx{}}~~\lstbox{\textSFx{}}~~\lstbox{\textSFx{}}~~\lstbox{\textSFx{}}~~\lstbox{\textSFx{}}~~\lstbox{\textSFx{}}~
\end{lstlisting}

\subsubsection*{How to fix it}

Declare the overriding method with \tokennolst{over yield}, and yield its values.

\errorcode{E4320}{generator \errhole{} cannot override the non generator method \errhole{}}
\label{sec:codes:e4320}

Raised during \textbf{validation}, from \texttt{ValidateErrorMessage::\allowbreak{}OVERRIDE\_\allowbreak{}NO\_\allowbreak{}GENERATOR\_\allowbreak{}BY\_\allowbreak{}GENERATOR} (\textit{src/\allowbreak{}ymirc/\allowbreak{}semantic/\allowbreak{}validator/\allowbreak{}errors.yr}).

\subsubsection*{What it means}

A generator, declared by \tokennolst{over yield}, tries to override an ordinary method, even one returning the handle of a generator, cf. \hyperref[sec:codes:e4319]{E4319}.

\subsubsection*{Example}

\textit{test\_\allowbreak{}resources/\allowbreak{}generators/\allowbreak{}test86.yr}:

%% check: skip
\begin{lstlisting}[style=coloredverbatimError]
// a generator cannot override a method returning the handle of a generator
yield gen ()-> i32 {
    yield 1;
}

class A {
    pub self () {}

    pub fn f (self)-> (yield-> i32) {
        gen ()
    }
}

class B over A {
    pub self () {}

    pub over yield f (self)-> i32 {
        yield 2;
    }
}

fn main () {
    let _ = B ();
}
\end{lstlisting}

\begin{lstlisting}[style=bashVerb, breaklines=true, escapechar=~]
Error[E4320] : generator test86::B::f ()-> (yield-> i32) cannot override the non generator method test86::A::f ()-> (yield-> i32)
 --> test_resources/generators/test86.yr:(17,20)
17  ~\lstbox{\textSFxi{}}~     pub over yield f (self)-> i32 {
    ~\lstbox{\textSFv{}}~                    ^
 --> test_resources/generators/test86.yr:(9,12)
 9  ~\lstbox{\textSFxi{}}~     pub fn f (self)-> (yield-> i32) {
    ~\lstbox{\textSFv{}}~            -
    ~\lstbox{\textSFxi{}}~
    = when validating test86::B -- (test_resources/generators/test86.yr:(14,7))
    ~\lstbox{\textSFii{}}~~\lstbox{\textSFx{}}~~\lstbox{\textSFx{}}~~\lstbox{\textSFx{}}~~\lstbox{\textSFx{}}~~\lstbox{\textSFx{}}~~\lstbox{\textSFx{}}~~\lstbox{\textSFx{}}~
Error[E4104] : symbol test86::B has errors
 --> test_resources/generators/test86.yr:(23,13)
23  ~\lstbox{\textSFxi{}}~     let _ = B ();
    ~\lstbox{\textSFv{}}~             ^
 --> test_resources/generators/test86.yr:(9,12)
 9  ~\lstbox{\textSFxi{}}~     pub fn f (self)-> (yield-> i32) {
    ~\lstbox{\textSFv{}}~            -
   ...
17  ~\lstbox{\textSFxi{}}~     pub over yield f (self)-> i32 {
    ~\lstbox{\textSFv{}}~                    -  error: generator test86::B::f ()-> (yield-> i32) cannot override the non generator method test86::A::f ()-> (yield-> i32)
    ~\lstbox{\textSFxi{}}~
    = when validating test86::B -- (test_resources/generators/test86.yr:(14,7))
    = when validating test86::main -- (test_resources/generators/test86.yr:(22,4))
    ~\lstbox{\textSFii{}}~~\lstbox{\textSFx{}}~~\lstbox{\textSFx{}}~~\lstbox{\textSFx{}}~~\lstbox{\textSFx{}}~~\lstbox{\textSFx{}}~~\lstbox{\textSFx{}}~~\lstbox{\textSFx{}}~
\end{lstlisting}

\subsubsection*{How to fix it}

Declare the overriding method as an ordinary method, with \tokennolst{over}.
